% Sesión 01: reconstrucción del contenido LaTeX original, sin PDF anexo

\begin{idea}[Cómo leer y utilizar el material]
Esta sección del libro contiene una guía de clase, dos laboratorios, un taller con
reporte, el suplemento y los ejercicios. Cada componente posee una fuente
LaTeX independiente y enlaces internos. \textbf{No se espera impartir todo
el volumen durante cuatro horas}: la clase introduce y demuestra; el resto
es trabajo autónomo o consulta. El solucionario está en el solucionario reservado al profesorado.
\end{idea}


% ===== Fuente modulos/00_proposito_y_ruta.tex =====

\subsection{Cómo usar esta sesión integrada}
\phantomsection\label{s01:mod:ruta}
\begin{idea}[Pregunta práctica que justifica toda la sesión]
Un equipo técnico redacta bitácoras en español. Queremos saber si un modelo entrenado
sobre lenguaje general representa y predice esos textos con la misma eficacia que uno
adaptado al vocabulario del dominio. Antes de decir que un modelo es ``mejor'', debemos
\emph{definir unidades, aprender probabilidades, evaluar documentos no vistos y medir incertidumbre}.
La pregunta nació de un problema observable, no de una ecuación presentada sin contexto.
\end{idea}
\begin{center}
\begin{tikzpicture}[>=Stealth, node distance=6mm and 8mm,
  paso/.style={draw=azul, rounded corners,fill=azulclaro,minimum height=12mm,
  text width=31mm,align=center,font=\small}]
\node[paso](uno){Informes en español};
\node[paso,right=of uno](dos){Caracteres, tokens y pares};
\node[paso,right=of dos](tres){Conteos y probabilidades};
\node[paso,below=of tres](cuatro){BPC y diferencias por documento};
\node[paso,left=of cuatro](cinco){Informe con incertidumbre};
\node[paso,left=of cinco](seis){Pregunta de investigación};
\draw[->,thick](uno)--(dos);\draw[->,thick](dos)--(tres);
\draw[->,thick](tres)--(cuatro);\draw[->,thick](cuatro)--(cinco);
\draw[->,thick](cinco)--(seis);
\end{tikzpicture}
\end{center}
\noindent\textbf{Dos preguntas, dos mediciones:} los bigramas de caracteres
cuantifican una perdida predictiva (BPC), mientras que el
\hyperref[s01:lab:ic10]{microexperimento BPE} cuantifica tokens por palabra
al cambiar el tokenizador. La segunda magnitud no mide por si sola calidad
probabilística; separar ambas evita una conclusión indebida sobre IC10.

\noindent\textbf{Qué contiene este libro.} Esta sección es la puerta de entrada a los materiales de S01, que también se
pueden consultar y ejecutar de forma independiente. Se puede leer únicamente la guía inicial, entrar después al
\hyperref[s01:lab:mat]{laboratorio matemático}, ejecutar el
\hyperref[s01:lab:comp]{laboratorio computacional} y su
\hyperref[s01:lab:ic10]{microexperimento de tokenización}, consultar el
\hyperref[s01:mod:sup]{suplemento} o resolver el
\hyperref[s01:mod:taller]{taller}. La \hyperref[s01:mod:invest]{ruta de investigación}
relaciona esos resultados con el documento de investigación y el objetivo de una
futura contribución científica en una revista de primer cuartil (Q1).

\begin{center}\small
\begin{tabular}{p{.23\textwidth}p{.65\textwidth}}
\toprule
\textbf{Componente} & \textbf{Función y evidencia que produce}\\\midrule
Guía y ejemplos & Construir el lenguaje matemático y comprender su utilidad.\\
Laboratorio matemático & Repetir cálculos a mano y demostrar identidades.\\
Laboratorio computacional & Generar resultados de ejemplo y archivos de auditoría en R y Python.\\
Taller + reporte & Interpretar cifras y producir un informe breve reproducible.\\
Suplemento & Justificar hipótesis, propiedades y límites de validez.\\
Investigación & Elegir dos áreas; identificar el camino a una pregunta novedosa y evaluable.\\\bottomrule
\end{tabular}
\end{center}
\begin{contraste}[Cuatro alcances que nunca se deben confundir]
\textbf{Cálculo demostrado:} se deduce de hipótesis declaradas.
\textbf{ILUSTRATIVO:} corpus o resultado fabricado para aprender.
\textbf{REAL REPRODUCIDO:} script ejecutado sobre un conjunto identificado
(que podría seguir siendo sintético).
\textbf{EVIDENCIA DE INVESTIGACIÓN:} requiere datos de dominio,
baselines fuertes, protocolo y validación externa. Los experimentos de esta
sesión pertenecen al segundo y tercer tipo; \textbf{no prueban la hipótesis de IC10}.
\end{contraste}

% ===== Fin fuente =====


% ===== Fuente modulos/10_guia.tex =====

\phantomsection\label{s01:mod:guia}

% ===== Fuente fuentes/apertura_visual.tex =====
\subsection{Parte I. ¿Cómo aprende una máquina a continuar una frase?}

\subsubsection*{1. La pregunta que abre el curso}
Imagine que una persona lee «El gato duerme sobre la...» y debe completar la frase. Tal vez proponga \emph{alfombra}, \emph{mesa} o \emph{cama}. Algunas continuaciones resultan más naturales que otras. Una máquina puede hacer una tarea análoga: asignar números entre cero y uno a las continuaciones posibles, de modo que su suma sea uno. Aún no necesitamos redes neuronales para entender esta operación.

\begin{center}
\begin{tikzpicture}[>=Stealth, every node/.style={font=\small}, token/.style={draw=azul,rounded corners=3pt,fill=azulclaro,minimum height=9mm,inner xsep=9pt}]
\node[token] (t) at (0,0) {El gato duerme sobre la};
\node[token,fill=verdeclaro,draw=verde] (a) at (6.1,1.0) {alfombra $\;0.50$};
\node[token,fill=verdeclaro,draw=verde] (b) at (6.1,0) {cama $\;0.30$};
\node[token,fill=verdeclaro,draw=verde] (c) at (6.1,-1.0) {mesa $\;0.20$};
\draw[->,thick] (t.east)--(a.west);\draw[->,thick] (t.east)--(b.west);\draw[->,thick] (t.east)--(c.west);
\end{tikzpicture}
\end{center}
\noindent\textit{Lectura de la figura.} Los valores se han inventado para ilustrar el mecanismo; no proceden de un LLM ni afirman que esas sean las únicas palabras posibles. La distribución condicional depende del contexto, de las unidades de representación y del modelo concreto.

\begin{idea}[Lo que aprenderemos a hacer en esta sesión]
Tomar una frase, separarla en unidades, contar cómo se combinan esas unidades, convertir los conteos en probabilidades y evaluar qué tan bien se ajustan las probabilidades a secuencias no utilizadas para el ajuste.
\end{idea}

\subsubsection*{2. ¿Para qué estudiar secuencias?}
No toda aplicación de un modelo de lenguaje consiste en completar frases. Las mismas ideas matemáticas sustentan otras preguntas:
\begin{center}
\begin{tabularx}{\textwidth}{>{\bfseries}p{3.0cm}X>{\raggedright\arraybackslash}p{3.4cm}}
\toprule Tarea & Pregunta reconocible & Objeto de entrada\\\midrule
Clasificación & ¿Esta reseña expresa aprobación o rechazo? & Palabras de la reseña\\
Traducción & ¿Qué secuencia en otro idioma corresponde a esta? & Secuencia de palabras\\
Documentos técnicos & ¿Qué evento se describe en este informe? & Texto especializado\\
Series de tiempo & ¿Cómo evolucionará la demanda mañana? & Mediciones ordenadas\\
Predicción de texto & ¿Qué símbolo sigue a lo observado? & Contexto escrito\\\bottomrule
\end{tabularx}
\end{center}
Una serie de tiempo no es texto; sin embargo, ambas tienen \emph{orden} y \emph{dependencia entre observaciones}. Más adelante estudiaremos modelos con memoria aprendida, atención y Transformers. Hoy construiremos las bases probabilísticas que hacen inteligibles sus objetivos de entrenamiento y evaluación. La intuición de comenzar por aplicaciones y ejemplos de secuencias está inspirada en el enfoque didáctico del capítulo 6 de \emph{Deep Learning with R} (primera edición); todas las figuras siguientes son originales.

\subsubsection*{3. Antes de las ecuaciones: ¿qué es una secuencia?}
Una \emph{secuencia} es una lista ordenada: cambiar el orden puede cambiar el significado o el fenómeno. «El perro persigue al gato» no expresa lo mismo que «El gato persigue al perro». En una serie de temperatura, $20,21,25$ no describe la misma evolución que $25,21,20$.

\begin{center}
\begin{tikzpicture}[every node/.style={font=\small},tok/.style={draw=azul,rounded corners=2pt,fill=azulclaro,minimum height=9mm,minimum width=18mm}, >=Stealth]
\node[anchor=west,font=\bfseries] at (-.2,1.4) {Texto};
\foreach \x/\w in {0/El,2/perro,4/persigue,6/al,8/gato} {\node[tok] at (\x,0.65){\w};}
\draw[->,azul,thick] (.8,.65)--(1.1,.65);\draw[->,azul,thick] (2.8,.65)--(3.1,.65);\draw[->,azul,thick] (4.8,.65)--(5.1,.65);\draw[->,azul,thick] (6.8,.65)--(7.1,.65);
\node[anchor=west,font=\bfseries] at (-.2,-.35) {Serie};
\foreach \x/\w in {0/20,2/21,4/25,6/24,8/23} {\node[tok,fill=verdeclaro,draw=verde] at (\x,-1.1){\w};}
\foreach \x in {0,2,4,6} {\draw[->,verde,thick] (\x+.85,-1.1)--(\x+1.1,-1.1);}
\end{tikzpicture}
\end{center}
\textbf{Pregunta de comprensión:} ¿qué se pierde si guardamos solo las frecuencias de las palabras de una oración, sin su orden? Se pierde, al menos, la información que permite distinguir quién realizó una acción sobre quién.

\subsubsection*{4. Primera decisión: ¿qué unidad llamaremos token?}
Un \emph{token} es una unidad de representación definida por una regla concreta. No debe confundirse con la palabra ordinaria: un token puede ser un carácter, una palabra, una pieza de palabra o un símbolo especial. La \emph{tokenización} transforma texto en una secuencia de tokens.

\begin{center}
\begin{tikzpicture}[>=Stealth, every node/.style={font=\small}, b/.style={draw=azul,rounded corners,fill=azulclaro,align=center,minimum height=11mm}]
\node[b,text width=40mm] (a) at (0,0){Texto original\\\texttt{El gato salta}};
\node[b,text width=35mm] (w) at (6.5,1.5){Por palabras\\\texttt{El | gato | salta}};
\node[b,text width=48mm] (c) at (6.5,0){Por caracteres\\\texttt{E | l | \_ | g | a | t | o | ...}};
\node[b,text width=48mm] (s) at (6.5,-1.5){Por subpalabras\\\texttt{El | gat | o | salt | a}};
\draw[->,thick] (a.east)--(w.west);\draw[->,thick] (a.east)--(c.west);\draw[->,thick] (a.east)--(s.west);
\end{tikzpicture}
\end{center}
\noindent Los ejemplos de segmentación son \textbf{didácticos}, no salidas atribuidas a un tokenizador comercial. El guion bajo representa un espacio exclusivamente en el dibujo; otras implementaciones emplean otras convenciones.

\subsubsection*{5. Del token al número: una representación comprensible}
Los algoritmos necesitan identificadores y, a menudo, vectores. Construyamos un vocabulario diminuto: $V=\{\text{el},\text{gato},\text{duerme},\text{fin}\}$, con identificadores $1,2,3,4$. La frase «el gato duerme» pasa a $(1,2,3)$.
\begin{center}
\begin{tikzpicture}[>=Stealth, every node/.style={font=\small},box/.style={draw=azul,fill=azulclaro,rounded corners,align=center,text width=36mm,minimum height=13mm}]
\node[box] (t) at (0,0){«el gato duerme»};
\node[box] (s) at (4.3,0){tokens\\el $\mid$ gato $\mid$ duerme};
\node[box] (i) at (8.6,0){identificadores\\$1,2,3$};
\draw[->,thick] (t)--(s);\draw[->,thick] (s)--(i);
\end{tikzpicture}
\end{center}
La \emph{codificación one-hot} representa el identificador 2 mediante el vector $(0,1,0,0)^\top$, donde solo una coordenada vale uno. Un \emph{embedding} asignaría al token un vector real aprendido, por ejemplo de dimensión $d$; sus valores no son necesariamente probabilidades ni tienen que sumar uno. En S02 estudiaremos formalmente esas representaciones. La separación conceptual es esencial: \textbf{tokenizar, identificar y representar vectorialmente son operaciones diferentes}.

\subsubsection*{6. ¿Qué veremos, y qué dejamos para más adelante?}
\begin{center}
\begin{tikzpicture}[>=Stealth, every node/.style={font=\small},s/.style={draw=azul,rounded corners,fill=azulclaro,text width=22mm,minimum height=11mm,align=center}]
\node[s] (a) at (0,0){Texto y\\tokens};
\node[s] (b) at (2.8,0){Pares y\\conteos};
\node[s] (c) at (5.6,0){Probabilidad\\condicional};
\node[s] (d) at (8.4,0){Pérdida y\\perplejidad};
\draw[->,thick] (a)--(b);\draw[->,thick] (b)--(c);\draw[->,thick] (c)--(d);
\node[font=\scriptsize,text width=10cm,align=center] at (4.2,-1.1){Más adelante: embeddings, redes neuronales, atención, Transformers, generación y aplicaciones};
\end{tikzpicture}
\end{center}
La S01 explica el mapa completo sin pretender enseñar todas las arquitecturas. Si el estudiante entiende por qué tiene sentido asignar probabilidades condicionales a tokens, estará preparado para comprender qué optimiza una red neuronal cuando aprende a predecirlos.
\clearpage

% ===== Fin fuente =====


% ===== Fuente fuentes/escalera_conteos.tex =====
\subsection{Parte II. De una frase a los bigramas y a las probabilidades}

\subsubsection*{7. Unigramas, bigramas y trigramas: primero con palabras}
Tomemos la oración \textbf{«el gato mira al perro»}. Separamos por espacios y obtenemos cinco tokens: \texttt{el}, \texttt{gato}, \texttt{mira}, \texttt{al}, \texttt{perro}. Los \emph{unigramas} son las unidades individuales; los \emph{bigramas} son pares \emph{consecutivos}; los \emph{trigramas} son grupos de tres unidades consecutivas.
\begin{center}
\begin{tikzpicture}[every node/.style={font=\small},t/.style={draw=azul,fill=azulclaro,minimum width=22mm,minimum height=9mm,rounded corners}]
\foreach \x/\w in {0/el,2.35/gato,4.7/mira,7.05/al,9.4/perro} {\node[t] at (\x,0){\w};}
\draw[very thick,verde](-1,-.72)--(3.35,-.72);\node[font=\scriptsize,verde] at (1.18,-1.0){bigrama: el gato};
\draw[very thick,vino](3.7,.68)--(10.5,.68);\node[font=\scriptsize,vino] at (7.1,1.0){trigrama: mira al perro};
\end{tikzpicture}
\end{center}
\begin{center}
\begin{tabular}{lll}\toprule
Unigramas (5) & Bigramas (4) & Trigramas (3)\\\midrule
el & el gato & el gato mira\\
gato & gato mira & gato mira al\\
mira & mira al & mira al perro\\
al & al perro & \\
perro & & \\\bottomrule
\end{tabular}
\end{center}
\noindent\textbf{Regla general:} con $T$ tokens y sin añadir marcadores, existen $T-n+1$ ventanas de longitud $n$ si $1\le n\le T$. La colección de ventanas conserva el orden \emph{dentro de cada ventana}, pero no necesariamente el orden entre ventanas si se usa una representación de bolsa.

\subsubsection*{8. ¿En qué difieren secuencia, conjunto y bolsa?}
Consideremos \texttt{gato gato perro}. Como secuencia es una lista ordenada de tres entradas. Como conjunto de tipos distintos, $\{\texttt{gato},\texttt{perro}\}$, pierde multiplicidades y orden. Una \emph{bolsa de palabras} conserva frecuencias: \texttt{gato}: 2, \texttt{perro}: 1, pero no el orden. Una \emph{bolsa de bigramas} cuenta ventanas de dos tokens conservando su identidad de par, aunque descarta sus posiciones. No son nociones intercambiables.
\begin{center}
\begin{tikzpicture}[font=\small, b/.style={draw=azul,rounded corners,minimum width=35mm,minimum height=10mm,align=center,fill=azulclaro}]
\node[b] at (0,0){Secuencia\\\texttt{gato,gato,perro}};
\node[b] at (4.3,0){Conjunto\\\texttt{\{gato,perro\}}};
\node[b] at (8.6,0){Bolsa\\\texttt{gato:2,perro:1}};
\end{tikzpicture}
\end{center}
\noindent\textbf{Pregunta:} ¿pueden «el gato muerde al perro» y «el perro muerde al gato» compartir una misma bolsa de palabras? Sí. Sus bigramas, en cambio, contienen información local diferente.

\subsubsection*{9. Ahora usemos tres símbolos: una cadena que podemos contar completa}
Pasamos a un ejemplo pequeño para hacer cuentas sin ordenador. La cadena \texttt{abacabbacb} tiene diez posiciones y, si no usamos marcadores de inicio y fin, nueve transiciones. La escribimos numerando sus posiciones:
\begin{center}
\begin{tikzpicture}[every node/.style={font=\small},token/.style={draw=azul,rounded corners=2pt,fill=azulclaro,minimum width=9mm,minimum height=9mm}]
\foreach \i/\v in {1/a,2/b,3/a,4/c,5/a,6/b,7/b,8/a,9/c,10/b} {
\node[token] at ({(\i-1)*1.03},0) (t\i) {\texttt{\v}};
\node[font=\scriptsize,gris] at ({(\i-1)*1.03},.75) {\i};
}
\foreach \i in {1,...,9} {\pgfmathtruncatemacro\j{\i+1} \draw[->,verde,thin] (t\i.south east) to[bend right=25] (t\j.south west);}
\end{tikzpicture}
\end{center}
\begin{center}
\begin{tabular}{c|ccccccccc}\toprule
Posición inicial &1&2&3&4&5&6&7&8&9\\\midrule
Bigrama & ab&ba&ac&ca&ab&bb&ba&ac&cb\\\bottomrule
\end{tabular}
\end{center}
En esta tabla, el primer símbolo del par es el \emph{contexto} y el segundo es la \emph{continuación}. Contemos transiciones sin ninguna fórmula:
\begin{center}
\begin{tabular}{c|ccc|c}\toprule
Contexto $\downarrow$ / siguiente $\to$ & $a$&$b$&$c$&Total\\\midrule
$a$&0&2&2&4\\
$b$&2&1&0&3\\
$c$&1&1&0&2\\\midrule
Total&3&4&2&9\\\bottomrule
\end{tabular}
\end{center}
\textbf{De la tabla a la matriz:} ahora sí tiene sentido escribir
\[
C=\begin{pmatrix}0&2&2\\2&1&0\\1&1&0\end{pmatrix},
\qquad\sum_{i,j}C_{ij}=9.
\]
La matriz no añade un concepto nuevo: \emph{es simplemente la misma tabla, sin rótulos}. Las filas y columnas se ordenan como $a,b,c$.

\subsubsection*{10. ¿Cómo se convierte un conteo en una predicción?}
Si hemos visto cuatro continuaciones después de $a$, dos fueron $b$ y dos fueron $c$. La frecuencia relativa de $b$ en ese contexto es $2/4=1/2$; para $c$ también $1/2$. Para $a$ fue $0/4=0$.
\begin{center}
\begin{tikzpicture}[font=\small,>=Stealth]
\node[draw,circle,fill=azulclaro,minimum size=12mm] (a) at (0,0) {$a$};
\node[draw,circle,fill=verdeclaro,minimum size=10mm] (b) at (4,1.2) {$b$};
\node[draw,circle,fill=verdeclaro,minimum size=10mm] (c) at (4,-1.2) {$c$};
\draw[->,thick,verde] (a)--node[above] {$2/4$}(b);
\draw[->,thick,verde] (a)--node[below] {$2/4$}(c);
\node[align=left] at (8,0){Contexto: $a$\\Siguiente: $b$ o $c$\\Probabilidades: $0.5,0.5$};
\end{tikzpicture}
\end{center}
Si dividimos cada fila de $C$ por su total, obtenemos una \emph{matriz de transición}:
\[
\widehat P=\begin{pmatrix}0&1/2&1/2\\2/3&1/3&0\\1/2&1/2&0\end{pmatrix}.
\]
Cada fila suma uno. La entrada $(i,j)$ se interpreta como probabilidad estimada del siguiente símbolo $j$ cuando el anterior fue $i$. A este estimador de probabilidades por frecuencias se le llama \emph{máxima verosimilitud} (MLE): derivaremos el nombre en la siguiente parte.

\subsubsection*{11. Un problema real para nuestro modelo pequeño: lo no observado}
La transición $b\to c$ no apareció en la cadena. El modelo por frecuencias le asigna probabilidad cero. Sin embargo, que no haya ocurrido en una muestra \emph{no demuestra} que sea imposible en textos nuevos.
\begin{center}
\begin{tikzpicture}[font=\small,>=Stealth]
\node[draw,rounded corners,fill=azulclaro,minimum height=10mm] (d) at (0,0){Entrenamiento: \texttt{abacabbacb}};
\node[draw,rounded corners,fill=vinoclaro,minimum height=10mm] (p) at (5.5,0){Prueba: \texttt{abca}};
\draw[->,thick] (d)--node[above,align=center,font=\scriptsize]{modelo de\\bigramas}(p);
\node[font=\small,vino] at (5.5,-.9){En prueba aparece $b\to c$: MLE da $0$};
\end{tikzpicture}
\end{center}
Esta dificultad motiva el \emph{suavizado}: reservar algo de probabilidad para resultados ausentes de la muestra. Con pseudoconteo $\alpha=1$ y tres símbolos:
\[
\widehat P_{\alpha=1}=\begin{pmatrix}1/7&3/7&3/7\\1/2&1/3&1/6\\2/5&2/5&1/5\end{pmatrix}.
\]
\textbf{Verificación a mano:} en la fila $a$ contamos $(0,2,2)$ y añadimos uno a cada columna: $(1,3,3)$. El denominador es $4+3=7$, no $6$. En la fila $b$, $(2,1,0)\mapsto(3,2,1)$ y dividimos por $6$.

\subsubsection*{12. ¿Qué predice exactamente un modelo de bigramas?}
Para \texttt{abca}, bajo el modelo suavizado, podemos evaluar las transiciones $a\to b$, $b\to c$, $c\to a$:
\[
\Pr(b\mid a)\Pr(c\mid b)\Pr(a\mid c)=\frac37\cdot\frac16\cdot\frac25=\frac1{35}.
\]
Este producto es la \emph{probabilidad condicional de las tres continuaciones dado el primer símbolo}. Para obtener la probabilidad conjunta de la secuencia hay que multiplicar también la probabilidad inicial. Por ejemplo, si $\Pr(X_1=a)=1/3$, entonces $\Pr(\texttt{abca})=1/105$. No debemos llamar \emph{conjunta} a $1/35$.

\begin{idea}[Lo que ya sabemos]
Sin usar cálculo avanzado hemos definido una unidad (token), formado ventanas (bigramas), construido un conteo, interpretado una probabilidad condicional y detectado por qué un modelo necesita tratar eventos ausentes. La matemática que sigue explica y generaliza exactamente estas operaciones.
\end{idea}
\clearpage

% ===== Fin fuente =====


% ===== Fuente fuentes/modelo_evaluacion.tex =====
\subsection{Parte III. Formalizar lo que acabamos de contar}

\subsubsection*{13. Del conteo al modelo de probabilidad}
Decir que una máquina asigna probabilidades obliga a precisar el objeto: sea $V$ un vocabulario finito de símbolos, $X_t\in V$ el símbolo en la posición $t$ y $X_{1:T}=(X_1,\ldots,X_T)$ una secuencia aleatoria. La notación $p(X_t=x\mid X_{1:t-1}=h)$ significa \emph{la probabilidad del siguiente símbolo $x$ dada una historia concreta $h$}. No es la probabilidad de que la máquina ``entienda'' una oración.

\begin{ejemplo}[Volvamos al dibujo de las continuaciones]
Si tras «El gato duerme sobre la» el vocabulario de juguete fuera $V=\{\text{alfombra},\text{cama},\text{mesa}\}$, una distribución válida es $(0.50,0.30,0.20)$. La probabilidad de «cama» en ese contexto sería $0.30$. Para un vocabulario real habría muchas más opciones, incluidas puntuación y símbolos especiales.
\end{ejemplo}

\subsubsection*{14. Regla de la cadena: de condicionales a oraciones enteras}
Sabemos por la definición de probabilidad condicional que $p(x_1,x_2)=p(x_1)p(x_2\mid x_1)$ cuando $p(x_1)>0$. Para tres posiciones:
\begin{align*}
p(x_1,x_2,x_3)&=p(x_1,x_2)\,p(x_3\mid x_1,x_2)\\
&=p(x_1)\,p(x_2\mid x_1)\,p(x_3\mid x_1,x_2).
\end{align*}
Al repetir el argumento obtenemos la \textbf{regla de la cadena}, donde todas las condicionales se definen en los prefijos con probabilidad positiva:
\[
\boxed{\displaystyle p(x_{1:T})=p(x_1)\prod_{t=2}^T p(x_t\mid x_{1:t-1}).}
\]
Es una identidad de probabilidad, no un supuesto de aprendizaje profundo. \emph{Cada factor utiliza potencialmente toda la historia disponible}; una red neuronal causal puede parametrizarlo, pero un modelo de conteos con ventana corta utiliza una aproximación.

\subsubsection*{15. Hipótesis de Markov: una simplificación, no una identidad}
En un modelo de bigramas aproximamos $p(x_t\mid x_{1:t-1})$ mediante $q(x_t\mid x_{t-1})$. Para trigramas usamos los dos símbolos anteriores. La idea de \emph{memoria de orden $m$} es retener solo los $m$ símbolos precedentes.
\begin{center}
\begin{tikzpicture}[>=Stealth,font=\small,b/.style={draw=azul,circle,fill=azulclaro,minimum size=10mm}]
\node[b] (x1) at (0,0){$X_1$};\node[b] (x2) at (3,0){$X_2$};\node[b] (x3) at (6,0){$X_3$};
\draw[->,thick] (x1)--(x2);\draw[->,thick] (x2)--(x3);
\draw[->,dashed,vino,thick] (x1) to[bend left=35] node[above,font=\scriptsize]{posible dependencia omitida} (x3);
\end{tikzpicture}
\end{center}
\begin{ejemplo}[El límite de mirar solo el símbolo anterior]
Sean $X_1,X_2$ monedas justas independientes y $X_3=X_1$. Entonces
\[
p(X_1=1,X_2=0,X_3=1)=\tfrac12\cdot\tfrac12\cdot1=\tfrac14.
\]
La aproximación de orden uno remplaza la tercera condicional por $p(X_3=1\mid X_2=0)=1/2$, y da $1/8$. La dependencia de $X_3$ con $X_1$ no queda explicada por $X_2$.
\end{ejemplo}

\subsubsection*{16. Por qué dividir por los conteos es máxima verosimilitud}
Para un contexto fijo $h$, supongamos que contamos $c_i$ continuaciones del símbolo $i$ y $N=\sum_i c_i>0$. Buscamos probabilidades $q_i$ no negativas que sumen uno. Bajo observaciones condicionalmente independientes dadas $h$ (modelo categórico), la verosimilitud, ignorando un factor combinatorio que no depende de $q$, es
\[
L(q)=\prod_{i=1}^{K} q_i^{c_i},\qquad
\ell(q)=\sum_{i:c_i>0} c_i\ln q_i.
\]
Para un máximo interior en el soporte observado, introducimos un multiplicador $\lambda$:
\[
\mathcal L(q,\lambda)=\sum_{i:c_i>0}c_i\ln q_i+\lambda(1-\sum_iq_i).
\]
Derivamos: $c_i/q_i-\lambda=0$, luego $q_i=c_i/\lambda$ para índices con conteo positivo. Sumando estas probabilidades se obtiene $\lambda=N$. La solución en todo el símplex, incluido su borde, es
\[
\boxed{\widehat q_i=\frac{c_i}{N}.}
\]
Los conteos cero reciben probabilidad cero. Esta fórmula explica por qué las frecuencias de la parte II no fueron una receta arbitraria. Si un contexto nunca aparece ($N=0$), la fórmula no define una fila: necesitamos una política explícita (retroceso, prior, etc.).

\subsubsection*{17. Suavizado Dirichlet: qué ganamos y qué sacrificamos}
El estimador anterior puede asignar ceros a eventos plausibles. En el modelo bayesiano utilizamos una distribución \emph{Dirichlet} sobre el vector de probabilidades: $q\sim\mathrm{Dir}(\alpha_1,\ldots,\alpha_K)$, donde $\alpha_i>0$. Como su densidad es proporcional a $\prod_iq_i^{\alpha_i-1}$ y la verosimilitud a $\prod_iq_i^{c_i}$, la posterior es
\[
q\mid c\sim\mathrm{Dir}(c_1+\alpha_1,\ldots,c_K+\alpha_K).
\]
La probabilidad predictiva posterior de una nueva observación, que coincide con la media de esa posterior, es
\[
\boxed{\Pr(X_{N+1}=i\mid c)=\frac{c_i+\alpha_i}{N+\sum_j\alpha_j}.}
\]
Con $\alpha_i=\alpha$, se conoce como suavizado aditivo. En el suplemento se deriva la esperanza de Dirichlet y se explican soporte, concavidad y dependencia del tamaño del vocabulario. \textbf{No se promete mejor rendimiento} por añadir pseudoconteos; el valor de $\alpha$ debe elegirse en datos de validación.

\subsubsection*{18. ¿Cómo sabemos si el modelo predice bien?}
Supongamos que una secuencia observada recibe probabilidad $q(x_{1:T})$. Una pérdida natural es el negativo del logaritmo de esa probabilidad: $-\ln q(x_{1:T})$. Si una probabilidad es pequeña, la pérdida es grande; si $q=0$ para un resultado observado, la pérdida es infinita.
\begin{center}
\begin{tabular}{c|c|c}\toprule
Probabilidad observada & $-\log_2q$ (bits) & Lectura\\\midrule
$1$ & $0$ & Sin sorpresa\\
$1/2$ & $1$ & Una decisión binaria ideal\\
$1/4$ & $2$ & Más sorpresa\\
$1/8$ & $3$ & Todavía más sorpresa\\\bottomrule
\end{tabular}
\end{center}
La \emph{pérdida logarítmica total} de una cadena se obtiene sumando contribuciones por posición. Si los modelos comparten convenciones de inicio, fin y tokenización, la media $T^{-1}\sum_t-\log q(x_t\mid x_{<t})$ permite compararlos sobre la misma secuencia. La guía distingue \emph{ejemplos construidos} de \emph{mediciones reales}; calcular una cifra no equivale a demostrar calidad predictiva fuera del ejemplo.

\subsubsection*{19. Entropía: incertidumbre antes de observar el resultado}
Sea $p=(p_1,\ldots,p_K)$ la distribución verdadera de un resultado aleatorio. La \emph{entropía en bits} es el promedio de la información $-\log_2p_i$ que aporta observar un resultado $i$:
\[
H_2(p)=-\sum_i p_i\log_2p_i,\qquad 0\log_2 0:=0.
\]
\begin{ejemplo}[Dos monedas que no sorprenden igual]
Una moneda justa tiene $p=(1/2,1/2)$ y entropía de un bit. Una moneda con $p=(3/4,1/4)$ tiene entropía
$-(3/4)\log_2(3/4)-(1/4)\log_2(1/4)\approx0.8113$ bits. La segunda es más predecible, en promedio, si conocemos su distribución.
\end{ejemplo}
El suplemento demuestra $0\le H_2(p)\le\log_2K$, incluyendo los casos de igualdad. Esta desigualdad permite interpretar la entropía antes de usarla como función de costo.

\subsubsection*{20. Entropía cruzada y KL: evaluar con probabilidades equivocadas}
Ahora distinguimos dos papeles: $p$ describe de dónde provienen los datos; $q$ es el modelo que les asigna probabilidades. La \emph{entropía cruzada} mide cuánto costaría codificar, en promedio, los datos de $p$ usando $q$:
\[
H_2(p,q)=-\sum_i p_i\log_2q_i.
\]
La diferencia con el costo ideal $H_2(p)$ se llama divergencia de Kullback--Leibler:
\[
D_{\rm KL}(p\|q)=\sum_i p_i\log_2\frac{p_i}{q_i},
\qquad H_2(p,q)=H_2(p)+D_{\rm KL}(p\|q).
\]
\begin{ejemplo}[Un ejemplo pequeño que se puede verificar]
Tomemos $p=(1/2,1/4,1/4)$ y $q=(1/4,1/2,1/4)$. Entonces
\begin{align*}
H_2(p)&=-\tfrac12(-1)-\tfrac14(-2)-\tfrac14(-2)=1.5,\\
H_2(p,q)&=-\tfrac12(-2)-\tfrac14(-1)-\tfrac14(-2)=1.75,\\
D_{\rm KL}(p\|q)&=1.75-1.5=0.25\ \text{bits}.
\end{align*}
El suplemento demuestra que KL es no negativa, que no es simétrica, que puede ser infinita y que no define una distancia métrica. La igualdad anterior no necesita ninguna red neuronal.
\end{ejemplo}

\subsubsection*{21. Perplejidad: deshacer la escala logarítmica}
Las pérdidas en bits son sumas convenientes, pero pueden ser poco intuitivas. Si una secuencia tiene pérdida promedio $L_2$ en bits por evento, definimos
\[
\boxed{\mathrm{PPL}=2^{L_2}}\qquad \text{o, en nats,}\qquad \boxed{\mathrm{PPL}=e^{L_e}}.
\]
\begin{center}
\begin{tikzpicture}[>=Stealth,font=\small,b/.style={draw=azul,rounded corners,fill=azulclaro,align=center,text width=33mm,minimum height=12mm}]
\node[b] (a) at (0,0){Probabilidades\\$1/2,1/4,1/8,1/2$};
\node[b] (b) at (4.1,0){Información\\$1,2,3,1$ bits};
\node[b] (c) at (8.2,0){Media $=1.75$\\PPL $=2^{1.75}$};
\draw[->,thick] (a)--(b);\draw[->,thick] (b)--(c);
\end{tikzpicture}
\end{center}
La perplejidad vale aproximadamente $3.364$. Es el equivalente uniforme de la pérdida media, no el número literal de palabras posibles ni una probabilidad de acierto. Si uno de los eventos observados recibe probabilidad cero, su PPL se vuelve infinita.

\subsubsection*{22. ¿Por qué no comparamos cualquier par de perplejidades?}
Los tokens de dos tokenizadores no representan necesariamente el mismo número de caracteres. Por eso, antes de comparar resultados debemos fijar el texto, la tokenización, la normalización, los marcadores y cómo se calcula la probabilidad de la secuencia completa.
\begin{ejemplo}[Dos números que engañan si se omite la unidad]
En el \emph{mismo texto de 100 caracteres}, supongamos que un modelo A utiliza 100 eventos con PPL $3.2$ y otro B utiliza 30 eventos con PPL $25$. Si ambos definen probabilidades coherentes de todo el texto y han tratado de forma compatible inicio y fin:
\[
\mathrm{BPC}_A=\frac{100\log_2(3.2)}{100}\approx1.678,
\qquad
\mathrm{BPC}_B=\frac{30\log_2(25)}{100}\approx1.393.
\]
B presenta una menor tasa bajo \emph{esas convenciones}, aunque tenga mayor PPL por token. Los números son ilustrativos; esta conversión no prueba una ventaja causal de su tokenizador. El análisis de condiciones y riesgos queda para el suplemento y el seminario.
\end{ejemplo}

\subsubsection*{23. ¿Por qué existen las subpalabras? Una puerta a BPE}
La tokenización por palabras puede producir vocabularios inmensos y palabras desconocidas; la tokenización por caracteres crea secuencias largas. Las subpalabras intentan encontrar unidades intermedias útiles. El algoritmo de \emph{byte pair encoding} (BPE), adaptado a la tokenización, comienza con unidades pequeñas y fusiona iterativamente pares vecinos frecuentes.
\begin{center}
\begin{tikzpicture}[>=Stealth,font=\small,b/.style={draw=azul,rounded corners,fill=azulclaro,align=center,text width=32mm,minimum height=11mm}]
\node[b] (a) at (0,0){\texttt{c | a | s | a}};
\node[b] (b) at (4.1,0){\texttt{c | as | a}};
\node[b] (c) at (8.2,0){\texttt{cas | a}};
\draw[->,thick] (a)--node[above,font=\scriptsize]{fusionar $a,s$}(b);
\draw[->,thick] (b)--node[above,font=\scriptsize]{fusionar $c,as$}(c);
\end{tikzpicture}
\end{center}
El dibujo muestra una \emph{posible} sucesión de fusiones, no el aprendizaje real de un tokenizador entrenado con corpus externos. El taller reproduce un ejemplo con palabras españolas y una regla determinista de desempate. BPE es un procedimiento voraz: no garantiza minimizar globalmente la longitud tras un número fijo de fusiones.

\subsubsection*{24. Repaso de las ideas y enlace con la investigación}
Hemos seguido la cadena \emph{texto $\to$ tokens $\to$ pares $\to$ conteos $\to$ probabilidades $\to$ pérdidas $\to$ comparación}. Cada paso exige convenciones que pueden registrarse, comprobarse y discutirse. Una pregunta legítima es si el texto técnico en español se segmenta de una manera distinta a la prosa general. Para responderla no basta observar dos ejemplos: hacen falta datos con procedencia, una hipótesis, métricas e incertidumbre. La parte IV introduce esa transición al trabajo científico.
\clearpage

% ===== Fin fuente =====


% ===== Fuente fuentes/casos_visuales.tex =====
\subsection{Parte IV. Dos casos integradores para conectar ideas y decisiones}

\subsubsection*{25. Caso A: del comentario de un usuario a una representación numérica}
Supongamos que una aplicación recibe tres frases: «servicio muy bueno», «servicio bueno» y «servicio muy lento». Antes de clasificar cada una como valoración positiva o negativa, debemos hacer explícitas varias decisiones: si distinguimos mayúsculas, si eliminamos tildes, si preservamos espacios y si trabajamos con palabras o subpalabras.

\begin{center}
\begin{tikzpicture}[>=Stealth,font=\small,b/.style={draw=azul,rounded corners,fill=azulclaro,minimum height=12mm,text width=31mm,align=center}]
\node[b] (a) at (0,0){Texto\\servicio muy bueno};
\node[b] (b) at (4.1,0){Tokens\\servicio $\mid$ muy $\mid$ bueno};
\node[b] (c) at (8.2,0){Códigos\\$1,2,3$};
\draw[->,thick] (a)--(b);\draw[->,thick] (b)--(c);
\end{tikzpicture}
\end{center}
Definamos, solo para este caso, $V=\{\texttt{servicio},\texttt{muy},\texttt{bueno},\texttt{lento}\}$. La bolsa de palabras de la primera frase es $(1,1,1,0)$ y la de la tercera $(1,1,0,1)$. Esas representaciones permiten comparar \emph{frecuencias}, pero no explican por sí mismas cómo se determina la polaridad de una opinión. Clasificar es una tarea distinta de predecir el siguiente token; un sistema moderno puede combinar representaciones y un predictor de clases.

\subsubsection*{26. Caso B: secuencia de eventos de un pozo, sin inventar datos reales}
Imagine una bitácora simbólica con eventos \texttt{N} (operación normal), \texttt{A} (alarma) y \texttt{R} (revisión): \texttt{N N A R N A R}. Aquí las letras son \emph{etiquetas ficticias} y no representan mediciones de un pozo real. La secuencia tiene siete posiciones y seis transiciones:
\[
N\to N,\quad N\to A,\quad A\to R,\quad R\to N,\quad N\to A,\quad A\to R.
\]
La matriz de conteos, con orden $N,A,R$, es
\[
C=\begin{pmatrix}1&2&0\\0&0&2\\1&0&0\end{pmatrix},\qquad
\widehat P=\begin{pmatrix}1/3&2/3&0\\0&0&1\\1&0&0\end{pmatrix}.
\]
La interpretación correcta es modesta: \emph{en esta secuencia ficticia}, una alarma fue seguida por una revisión en las dos ocasiones observadas. No significa que una alarma en un sistema industrial real siempre produzca una revisión. Para esa afirmación se requerirían fuentes, datos, diseño y validación.

\begin{idea}[Texto y series de tiempo]
Una secuencia de eventos discretos puede estudiarse con una cadena de Markov. Una serie temporal de presión o temperatura contiene valores numéricos y puede requerir un modelo distinto. La conexión entre ambos cursos es estructural (dependencias y pronóstico), no la afirmación de que todas las secuencias se comportan como palabras.
\end{idea}

\subsubsection*{27. Un cálculo adicional con entropía y perplejidad}
Para el contexto $A$ de la bitácora de juguete, el estimador MLE asigna $(0,0,1)$, y la entropía de su continuación es cero: según ese modelo, no existe incertidumbre. Con suavizado aditivo $\alpha=1$, los conteos $(0,0,2)$ producen probabilidades $(1/5,1/5,3/5)$. Su entropía en bits es
\[
H=-2\cdot\tfrac15\log_2(\tfrac15)-\tfrac35\log_2(\tfrac35)\approx1.371.
\]
\textbf{¿Contradicción?} No. La entropía mide incertidumbre \emph{según una distribución específica}; al cambiar el modelo, cambia la entropía del modelo. No hemos demostrado una variación de la incertidumbre real del fenómeno.

Un experimento evaluativo debe separar \emph{modelo ajustado} y \emph{datos de prueba}. Si una nueva transición observada fuera $A\to N$, bajo MLE produciría pérdida infinita, mientras que bajo el suavizado anterior costaría $-\log_2(1/5)=\log_2 5\approx2.322$ bits. La decisión de usar suavizado influye en la generalización y en la función de pérdida.

\subsubsection*{28. Cómo evitar conclusiónes precipitadas}
Considere cuatro afirmaciones que se pueden discutir antes de emprender un proyecto:
\begin{enumerate}
\item «Una palabra aparece más veces en el corpus». \emph{Descriptiva}, siempre que se defina el corpus y la normalización.
\item «Un tokenizador produce menos tokens por palabra». \emph{Medición comparativa} que exige textos y vocabularios fijados.
\item «Un modelo produce menor perplejidad». \emph{Resultado dependiente del protocolo}: se deben declarar tokenización y unidad.
\item «Este tokenizador produce mejores respuestas». \emph{Hipótesis causal o predictiva más fuerte}, que requiere controlar modelo, entrenamiento, datos y métricas de tarea.
\end{enumerate}
Estas diferencias separan una demostración pedagógica de una investigación reproducible. Conducen directamente a la última hora del seminario y al hito de dos áreas de interés.
\clearpage

% ===== Fin fuente =====


% ===== Fuente fuentes/bpe_paso_a_paso.tex =====
\subsection{Parte V. BPE paso a paso y buenas prácticas de comparación}

\subsubsection*{29. Cinco palabras, una regla de fusión}
La técnica de BPE puede ilustrarse sin emplear bibliotecas de aprendizaje profundo. Usaremos las palabras \emph{casa} (5 apariciones), \emph{casas} (2), \emph{cosa} (6), \emph{cosas} (3) y \emph{caso} (4). Cada palabra se representa inicialmente como caracteres y un marcador de fin $\langle/\mathrm{w}\rangle$. La frecuencia de una palabra pondera todas sus parejas adyacentes. El objetivo de este ejercicio no es crear un tokenizador listo para producción, sino entender qué se modifica al aprender unidades.

\begin{center}
\begin{tabular}{lcl}\toprule
Palabra & Frecuencia & Segmentación inicial\\\midrule
casa &5& c $\mid$ a $\mid$ s $\mid$ a $\mid$ $\langle/\mathrm{w}\rangle$\\
casas &2& c $\mid$ a $\mid$ s $\mid$ a $\mid$ s $\mid$ $\langle/\mathrm{w}\rangle$\\
cosa &6& c $\mid$ o $\mid$ s $\mid$ a $\mid$ $\langle/\mathrm{w}\rangle$\\
cosas &3& c $\mid$ o $\mid$ s $\mid$ a $\mid$ s $\mid$ $\langle/\mathrm{w}\rangle$\\
caso &4& c $\mid$ a $\mid$ s $\mid$ o $\mid$ $\langle/\mathrm{w}\rangle$\\\bottomrule
\end{tabular}
\end{center}
\noindent En este corpus de juguete, el par de caracteres $(a,s)$ aparece ponderadamente $5+2\cdot2+3+4=16$ veces (en \emph{casas} aparece dos veces). El par $(s,a)$ también suma 16: hay un empate que resolvemos con la regla lexicográfica declarada. Se elige $(a,s)$ y se fusiona en una unidad \texttt{as}. En \emph{casa} pasamos de \texttt{c|a|s|a|fin} a \texttt{c|as|a|fin}. El cambio no altera el texto original; altera la forma de \emph{segmentarlo}.

\subsubsection*{30. Cómo leer una tabla de fusiones}
Las seis primeras fusiones, bajo conteo ponderado, empates lexicográficos y sustitución de pares no solapados, resultan así:
\begin{center}
\begin{tabular}{c l r r}\toprule
Paso & Par fusionado & Frecuencia & Unidades después\\\midrule
0 & -- & -- &105\\
1 & $(a,s)$ &16&89\\
2 & $(a,\langle/\mathrm{w}\rangle)$ &11&78\\
3 & $(c,as)$ &11&67\\
4 & $(c,o)$ &9&58\\
5 & $(co,s)$ &9&49\\
6 & $(cos,a\langle/\mathrm{w}\rangle)$ &6&43\\\bottomrule
\end{tabular}
\end{center}
\noindent Cada resta $105-89=16$, $89-78=11$, etc. es verificable en el módulo R/Python. No debemos suponer sin comprobar que \emph{frecuencia de un par} y \emph{sustituciones efectivas} siempre coinciden: en \texttt{aaaaa}, el par $(a,a)$ aparece en cuatro ventanas superpuestas, pero una pasada de fusión no solapada puede efectuar solamente dos reemplazos.

\subsubsection*{31. Una tabla de preguntas para discutir con la clase}
\begin{center}
\begin{tabularx}{\textwidth}{p{55mm}X}\toprule
Pregunta & Idea matemática o metodológica\\\midrule
¿BPE siempre disminuye el número de unidades? & Una fusión aplicada sustituye dos unidades por una; las reducciones efectivas dependen de los solapamientos.\\
¿BPE produce mejores predicciones? & No se deduce de reducir longitud: es necesario evaluar modelos probabilísticos en prueba.\\
¿Podemos comparar PPL entre segmentaciónes? & No directamente por token; fijar texto, unidades y probabilidades de la cadena completa.\\
¿El corpus técnico permite generalizar al español técnico? & Solo con muestreo adecuado, procedencia y dominio especificados.\\\bottomrule
\end{tabularx}
\end{center}
\textbf{Puente hacia el proyecto IC10.} Podemos investigar si ciertos términos técnicos ocupan más tokens por palabra que la prosa general. Esto ya es una medición falsable y puede ampliarse a BPC con modelos adecuadamente controlados. La tabla del catálogo identifica D01 y D02 como corpus generales candidatos; ninguna cifra empírica sobre esos corpus se afirma en esta guía.
\begin{idea}[Tres capas: prompt, contexto y arnés]
En un modelo generativo, el \emph{prompt} aporta la entrada inicial; el \emph{contexto} es la historia efectiva disponible, según la ventana y las reglas de truncamiento; el \emph{arnés} ejecuta y registra el procedimiento (muestreo, parada, validación, semillas y herramientas). Que los tres influyan en la salida no los convierte en la misma operación. Su estudio sistemático se retomará en sesiones posteriores.
\end{idea}
\clearpage

% ===== Fin fuente =====


% ===== Fuente fuentes/glosario_consulta.tex =====
\subsection{Vocabulario de consulta: diez ideas en una página}

Los términos de esta tabla se han construido en la guía. Se reúnen aquí para que el estudiante pueda volver a sus definiciones sin perder el hilo de la lectura.
\begin{center}
\begin{tabularx}{\textwidth}{>{\bfseries}p{37mm}X}\toprule
Término & Qué significa en esta sesión \\\midrule
Secuencia & Lista ordenada de observaciones; el orden puede afectar su significado o distribución.\\
Token & Unidad que produce una regla de segmentación; puede ser palabra, carácter o subpalabra.\\
Tokenización & Procedimiento que convierte texto en tokens, con convenciones explícitas.\\
N-grama & Ventana de $n$ unidades consecutivas; un bigrama contiene dos.\\
Bolsa de palabras & Recuento de tokens que ignora su posición en la secuencia.\\
Contexto & Información ya disponible al estimar una probabilidad de continuación.\\
MLE & Elección de parámetros que maximiza la probabilidad asignada a la muestra de entrenamiento.\\
Suavizado & Procedimiento que modifica probabilidades estimadas para evitar asignar cero por ausencia muestral.\\
Entropía cruzada & Pérdida logarítmica media de datos de $p$ al usar probabilidades de $q$.\\
Perplejidad & Exponencial de la pérdida logarítmica media por evento, con unidad de evaluación declarada.\\\bottomrule
\end{tabularx}
\end{center}
\begin{idea}[¿Qué viene después?]
El token es la unidad; el vector es una representación; el modelo produce una distribución condicional; la pérdida permite evaluar esa distribución; la investigación pregunta si diferencias observadas se sostienen más allá de un ejemplo elegido a mano.
\end{idea}
\clearpage

% ===== Fin fuente =====


% ===== Fin fuente =====


% ===== Fuente modulos/20_laboratorio_matematico.tex =====

\phantomsection\label{s01:lab:mat}

% ===== Fuente fuentes/ejemplos_verificados.tex =====
\subsection{Laboratorio matemático de pizarra: estimar y comprobar}


\subsubsection*{Un experimento completo desde cero: cuatro símbolos en prueba}

Fijemos el alfabeto $V=\{a,b,c\}$, una cadena de entrenamiento
$\texttt{abacabbacb}$ y la cadena reservada $\texttt{abca}$.
Ambas son \textbf{ILUSTRATIVAS}: se construyeron para enseñar la cuenta, no provienen de
un corpus observado. No obstante, sus resultados numéricos están regenerados por
\texttt{codigo/s01.py} y contrastados en \texttt{verificaciones/}.

Construimos la matriz de transiciones $C_{ij}=\#\{t:x_t=i,x_{t+1}=j\}$:
\[
C=\begin{array}{c|rrr|r}&a&b&c&N_i\\\hline
 a&0&2&2&4\\ b&2&1&0&3\\ c&1&1&0&2\end{array}.
\]
Hay $10-1=9$ transiciones, no diez: el símbolo inicial no tiene predecesor.
La ecuación de normalización para $\alpha=1$ es
\[
 P_{ij}=\frac{C_{ij}+1}{N_i+3}
 =\begin{pmatrix}1/7&3/7&3/7\\1/2&1/3&1/6\\2/5&2/5&1/5\end{pmatrix}.
\]
El denominador de la primera fila es $7$, no $6$; el de la tercera es $5$.
La comprobación de suma por fila es elemental pero obligatoria:
$1/7+3/7+3/7=1$, $1/2+1/3+1/6=1$ y $2/5+2/5+1/5=1$.
La probabilidad \emph{condicional de las tres transiciones de prueba}, dado $X_1=a$, es
\[
 P(b\mid a)P(c\mid b)P(a\mid c)
 =\frac37\cdot\frac16\cdot\frac25=\frac1{35}.
\]
No es todavía la probabilidad de la cadena completa. Para ello hay que precisar
la distribución inicial $\pi$: si elegimos $\pi=(1/3,1/3,1/3)$,
\[
 q(\texttt{abca})=\pi(a)\,P(b\mid a)P(c\mid b)P(a\mid c)=\frac1{105}.
\]
Por tanto, la longitud ideal de código es $\log_2(105)\approx6.714246$ bits,
y, con cuatro eventos de predicción, la tasa es
$\log_2(105)/4\approx1.678561$ bits por carácter y la perplejidad
$105^{1/4}\approx3.201086$.
Si en lugar de cuatro eventos se evalúan solo las tres transiciones,
la pérdida media pasa a $\log_2(35)/3$; \emph{ambas cifras son válidas,
pero corresponden a experimentos condicionales diferentes}.

\begin{ejemplo}[Un contraste que detecta un modelo incorrecto]
El estimador no suavizado asigna $\widehat P(c\mid b)=0$.
Entonces $q(\texttt{abca})=0$ para cualquier $\pi(a)>0$ y la pérdida
es $+\infty$. Sustituir el cero únicamente al mirar la prueba sería
modificar el modelo después de observarla: hay que definir el suavizado y
seleccionar su intensidad en validación.
\end{ejemplo}

\begin{figure}[H]\centering
\begin{tikzpicture}[>=Stealth,font=\small, every node/.style={align=center}]
\node[draw=azul,fill=azulclaro,rounded corners,minimum width=3cm,minimum height=1cm] (t) at (0,0) {\textbf{Entrenamiento}\\10 símbolos};
\node[draw=azul,fill=azulclaro,rounded corners,minimum width=3cm,minimum height=1cm] (c) at (4,0) {\textbf{Conteos}\\9 transiciones};
\node[draw=verde,fill=verdeclaro,rounded corners,minimum width=3cm,minimum height=1cm] (p) at (8,0) {\textbf{Suavizado}\\filas de $P$};
\node[draw=ocre,fill=ocreclaro,rounded corners,minimum width=3cm,minimum height=1cm] (e) at (8,-2) {\textbf{Prueba}\\$\texttt{abca}$};
\node[draw=azul,fill=azulclaro,rounded corners,minimum width=3cm,minimum height=1cm] (r) at (2,-2) {\textbf{Evaluación}\\$\log_2 105$ bits};
\draw[->,thick] (t)--(c);\draw[->,thick] (c)--(p);\draw[->,thick] (p)--(e);\draw[->,thick] (e)--(r);
\end{tikzpicture}
\caption{La prueba solo recibe un modelo previamente fijado. El flujo contrario constituye fuga de información.}
\end{figure}

\subsubsection*{Regla de la cadena frente a memoria finita: una cuenta contrastable}

Tomemos $X_1$ y $X_2$ Bernoulli$(1/2)$ independientes, y definamos
$X_3=X_1$ casi seguramente. Se trata de una distribución conjunta legítima,
aunque no sea markoviana de orden uno. Para $x=(1,0,1)$,
\[
 p(1,0,1)=p(X_1=1)\,p(X_2=0\mid X_1=1)\,
 p(X_3=1\mid X_1=1,X_2=0)=\frac12\frac12\cdot1=\frac14.
\]
Un modelo construido con los \emph{verdaderos} condicionales de orden uno
utilizaría $p(X_3=1\mid X_2=0)=1/2$ y asignaría $1/8$ a esa cadena.
Esta diferencia no es error algebraico: proviene de descartar información.

\begin{figure}[H]\centering
\begin{tikzpicture}[>=Stealth,font=\small,
 v/.style={draw=azul,fill=azulclaro,circle,minimum size=10mm}]
\node[v] (x1) at (0,0) {$X_1$};\node[v] (x2) at (4,0) {$X_2$};\node[v] (x3) at (8,0) {$X_3$};
\draw[->,thick] (x1)--(x2);\draw[->,thick] (x2)--(x3);
\draw[->,thick,dashed,color=vino] (x1) to[bend left=35] node[above] {memoria relevante} (x3);
\node at (4,-1.0) {Un bigrama descarta la dependencia discontinua $X_1\rightarrow X_3$.};
\end{tikzpicture}
\caption{Diagrama explicativo de una dependencia que no cabe en un único símbolo de memoria.}
\end{figure}

\subsubsection*{Suavizado como regularización: el papel del tamaño del vocabulario}

Sea $K$ el tamaño del alfabeto y $\alpha>0$. El peso total de los
pseudoconteos es $K\alpha$, de modo que
\[
\frac{c_i+\alpha}{N+K\alpha}
=\underbrace{\frac{N}{N+K\alpha}}_{\lambda}\frac{c_i}{N}
 +\underbrace{\frac{K\alpha}{N+K\alpha}}_{1-\lambda}\frac1K.
\]
Esta identidad describe el suavizado como mezcla convexa entre el estimador
empírico y la distribución uniforme. Es válida para $N>0$ y $\alpha>0$.
Para $N=10$, $K=5$ y $\alpha=1$, $\lambda=2/3$; si se aumenta
$K$ a 100, manteniendo $N$ y $\alpha$, $\lambda=1/11$. La masa que
se introduce por suavizado depende fuertemente del vocabulario.

\begin{ejemplo}[Una decisión que exige validación]
En un contexto con conteos $(6,3,1,0,0)$, $K=5$ y $\alpha=1$,
se obtiene $q=(7,4,2,1,1)/15$. El símbolo ausente recibe $1/15$.
Con $\alpha=1/2$, el denominador es $12.5$ y recibe $0.5/12.5=1/25$.
Ninguna de las dos opciones es automáticamente superior: la pérdida en
validación sirve para decidir antes de examinar prueba.
\end{ejemplo}

\subsubsection*{Perplejidad: tres lecturas compatibles y una lectura incorrecta}

Fijadas las probabilidades $r_t=q(x_t\mid h_t)>0$ de $T$ eventos observados,
\[
\operatorname{PPL}
=\exp\left(-\frac1T\sum_{t=1}^T\ln r_t\right)
=\left(\prod_{t=1}^T\frac1{r_t}\right)^{1/T}.
\]
En consecuencia, la PPL es la \emph{media geométrica de las inversas}
de las probabilidades asignadas a los eventos. Si todas son iguales a
$1/K$, la PPL es $K$; para probabilidades variables, la PPL no es
el número real de símbolos candidatos ni la inversa de la exactitud.

\begin{ejemplo}[Producto, bits y exponencial: cálculo completo]
Para probabilidades observadas $(1/2,1/4,1/8,1/2)$,
\[
-\sum_{t=1}^4\log_2r_t=1+2+3+1=7\text{ bits},
\qquad\operatorname{PPL}=2^{7/4}\approx3.363586.
\]
Si la tercera probabilidad se reduce de $1/8$ a $1/16$ sin alterar las demás,
la suma aumenta a $8$ bits y la PPL a $2^2=4$. El cambio puede
explicarse multiplicativamente: se duplicó el código ideal de ese evento.
\end{ejemplo}

\subsubsection*{Comparar tasas, no solo nombres de métricas}

Dos modelos evalúan exactamente la misma cadena de $100$ caracteres,
sin pérdida de información al tokenizar. El modelo $A$ predice
$100$ caracteres con PPL $3.2$, mientras que $B$ predice $30$ subpalabras
con PPL $25$. Sus longitudes ideales son
\[
B_A=100\log_2(3.2)\approx167.8072\text{ bits},\quad
B_B=30\log_2(25)\approx139.3157\text{ bits}.
\]
Si el texto de referencia y el cálculo probabilístico sobre \emph{todo} el texto
son comparables, las tasas respectivas son $1.6781$ y $1.3932$ bits por
carácter. En este ejemplo el orden por PPL es inverso al orden por tasa.

\begin{nota}[Limitación esencial]
La misma unidad de longitud no convierte esta comparación en un efecto
\emph{causal} de la tokenización: también pueden variar datos de entrenamiento,
normalización, parametrización, contexto y presupuesto computacional. En particular,
no se identifica con BPC la calidad semántica de un modelo.
\end{nota}
Tampoco debe contarse el mismo dato dos veces: si una cadena Unicode se
normaliza antes del modelado, el número de caracteres del denominador debe
reflejar una convención fija y documentada. Para comparar código sin pérdida
se requiere poder reconstruir exactamente la cadena original.

\subsubsection*{De muestras a poblaciones: qué se demuestra realmente}

La pérdida observada sobre $T$ eventos es
$\widehat R_T(q)=T^{-1}\sum_{t=1}^T-\ln q(x_t\mid h_t)$.
En un problema categórico \emph{i.i.d.} con observaciones de distribución
$p$ y soporte compatible, la esperanza de $-\ln q(X)$ es $H(p,q)$;
la ley de los grandes números justifica convergencia para $q$ fijo
si la variable es integrable. En una secuencia dependiente la convergencia
requiere hipótesis apropiadas (por ejemplo, ergodicidad e integrabilidad),
y en modelos elegidos usando los mismos datos aparece un problema adicional
de sobreajuste.

Por otro lado, para una muestra fija, minimizar $\widehat R_T(q_\theta)$
es exactamente maximizar la verosimilitud correspondiente, porque
\[
\arg\min_\theta\widehat R_T(q_\theta)
=\arg\min_\theta\left(-\frac1T\ln q_\theta(x_{1:T})\right)
=\arg\max_\theta q_\theta(x_{1:T}).
\]
Esto no prueba que $q_\theta=p$, ni establece por sí solo generalización.
El paso de lo observado a una afirmación científica exige particiones,
incertidumbre y una pregunta experimental bien definida.

% ===== Fin fuente =====


% ===== Fuente fuentes/derivaciones_de_pizarra.tex =====

\subsection{Derivaciones complementarias y experimentos de pizarra}


\subsubsection*{De una frase a una variable aleatoria}

Antes de introducir redes neuronales separamos tres operaciones distintas:
(i) elegir una representación discreta del texto, (ii) asignar probabilidades a
secuencias en esa representación y (iii) generar símbolos condicionando en un
prefijo. Por ejemplo, la cadena ``pozo estable'' puede representarse por caracteres,
palabras o subpalabras; las probabilidades sobre los elementos de cada representación
no tienen por qué coincidir. Una representación admite un modelo autorregresivo
si permite definir $q(x_t\mid x_{<t})$ y su producto a lo largo de la secuencia.

\begin{figure}[H]\centering
\begin{tikzpicture}[font=\small,box/.style={draw=azul,rounded corners=2pt,fill=azulclaro,
minimum height=9mm,align=center,text width=3cm},>=Stealth]
\node[box] (a) at (0,0) {Texto original\\\texttt{pozo estable}};
\node[box] (b) at (4.1,0) {Tokenización\\$x_1,\ldots,x_T$};
\node[box] (c) at (8.2,0) {Probabilidad\\$\prod_tq(x_t\mid x_{<t})$};
\draw[->,thick] (a)--(b);
\draw[->,thick] (b)--(c);
\node[draw=gris,dashed,rounded corners,text width=7.7cm,align=center] at (4.1,-1.5)
{Medición: $-\log_2 q(x_{1:T})$ bits; fijar fronteras y unidad del texto};
\end{tikzpicture}
\caption{Tres decisiones separadas: representación, asignación de probabilidad y medición.}
\end{figure}

\subsubsection*{Regla de la cadena con números: dependencia que el bigrama pierde}

Sea $X_1\sim\mathrm{Bernoulli}(1/2)$, $X_2\sim\mathrm{Bernoulli}(1/2)$ independiente
y $X_3=X_1$ con probabilidad uno. La regla de la cadena es exacta:
$p(x_1,x_2,x_3)=p(x_1)p(x_2\mid x_1)p(x_3\mid x_1,x_2)$.
La probabilidad de $(1,0,1)$ es $1/2\cdot1/2\cdot1=1/4$. Sin embargo,
$p(X_3=1\mid X_2=0)=1/2$; al imponer el supuesto de bigramas se obtiene
$1/2\cdot1/2\cdot1/2=1/8$. La diferencia se debe a la memoria perdida, no
a una falla de la regla de la cadena. En esta distribución el pasado remoto
$X_1$ es indispensable para predecir $X_3$.

\subsubsection*{La estimación por máxima verosimilitud cuando hay ceros}

Para un contexto $h$, sean $c_1,\ldots,c_K$ los conteos y $N=\sum_i c_i>0$.
El objetivo $\ell(q)=\sum_{i:c_i>0}c_i\ln q_i$ se maximiza sobre el símplex
cerrado $q_i\geq0$, $\sum_iq_i=1$. Sobre el soporte observado, las condiciones
de Lagrange dan $c_i/q_i=\lambda$. Por ello $q_i=c_i/N$, y por el principio
de la desigualdad de Gibbs este valor maximiza globalmente la log-verosimilitud.
Para un $c_i=0$, la solución asigna $q_i=0$; una solución estrictamente positiva
para todas las categorías no alcanza el máximo, solo puede aproximarlo.

\begin{center}
\begin{tabular}{lrrrr}\toprule
Contexto $h$ & $c(a)$ & $c(b)$ & $c(c)$ & $N_h$\\\midrule
$a$ & 0 & 2 & 2 & 4\\
$b$ & 2 & 1 & 0 & 3\\
$c$ & 1 & 1 & 0 & 2\\\bottomrule
\end{tabular}
\end{center}
Para el contexto $b$, $\hat q(c\mid b)=0$. Con $\alpha=1/2$ y $K=3$,
\[
q_\alpha(\cdot\mid b)=\biggl(\frac{2.5}{4.5},\frac{1.5}{4.5},
\frac{0.5}{4.5}\biggr)=\left(\frac59,\frac13,\frac19\right).
\]
La masa para el símbolo no visto es $1/9$; para $\alpha=1$ sería $1/6$.
Esto pone de relieve que el parámetro de suavizado cambia predicciones,
por lo que se selecciona en validación, nunca en prueba.

\subsubsection*{De Dirichlet a predicción posterior}

Sea $\mathbf q\sim\operatorname{Dirichlet}(\alpha_1,\ldots,\alpha_K)$.
Su densidad (hasta constante) es $\prod_iq_i^{\alpha_i-1}$.
Dada una muestra con conteos $\mathbf c$, la verosimilitud es proporcional
a $\prod_iq_i^{c_i}$, de modo que la posterior tiene densidad
$\prod_iq_i^{c_i+\alpha_i-1}$; por identificación,
$\mathbf q\mid \mathbf c\sim\operatorname{Dirichlet}(\mathbf c+\boldsymbol\alpha)$.
La distribución predictiva para una nueva observación es
\[
P(X_{N+1}=i\mid\mathbf c)=\int q_i\,p(\mathbf q\mid\mathbf c)\,d\mathbf q
=\frac{c_i+\alpha_i}{N+\sum_j\alpha_j}.
\]
La integral es una esperanza posterior, no una moda posterior ni un estimador
frecuentista sin regularización. Este punto facilita comprender que un modelo
puede conservar incertidumbre sobre parámetros a partir de pocos datos.

\subsubsection*{Entropía, cruzada, divergencia y codificación}

Para $p=(3/4,1/4)$ y $q=(1/2,1/2)$, usando logaritmos en base dos,
\[
H(p)=-\tfrac34\log_2\tfrac34-\tfrac14\log_2\tfrac14\simeq0.8113,
\qquad H(p,q)=1.
\]
Luego $D_{\mathrm{KL}}(p\Vert q)=H(p,q)-H(p)\simeq0.1887$ bits.
Invirtiendo el orden se obtiene
$D_{\mathrm{KL}}(q\Vert p)=\frac12\log_2(2/3)+\frac12\log_2 2
\simeq0.2075$ bits: KL no es simétrica.
El argumento de Gibbs establece $D_{\mathrm{KL}}(p\Vert q)\geq0$;
si $q=0$ en un punto con $p>0$, se entiende que la divergencia es infinita.
La interpretación de codificación es operacional: $-\log_2 q(x)$ mide
longitud idealizada de un código probabilístico para el evento $x$.

\subsubsection*{Tres modelos, una cadena: por qué la unidad importa}

Considérese un texto de $100$ caracteres Unicode medidos como puntos de
código, sin normalización posterior. Supónganse dos distribuciones válidas
sobre sus secuencias de tokens, con codificación reversible del texto:
\begin{center}
\begin{tabular}{lrrrr}\toprule
Modelo & Tokens & PPL/token & Bits totales & BPC\\\midrule
A (caracteres) &100& $3.2$ & $100\log_2(3.2)\simeq167.8$ & $1.678$\\
B (subpalabras)&30& $25$ & $30\log_2(25)\simeq139.3$ & $1.393$\\\bottomrule
\end{tabular}
\end{center}
Ambos juegos de probabilidades son \emph{ilustrativos}. Para comparar tasas
sobre un mismo texto debe incluirse toda la probabilidad del evento textual,
incluidos marcadores de inicio y fin cuando correspondan, y garantizar que
ambas representaciones reconstruyan el texto sin pérdida. Si un tokenizador
normaliza y destruye caracteres, el BPC calculado sobre texto original
puede dejar de representar una tasa comparable; es necesario especificar
la transformación o evaluarla explícitamente.

\subsubsection*{Ventana de contexto y mecanismo de generación}

En el generador autorregresivo, el \emph{prompt} proporciona el prefijo;
el contexto es la historia efectivamente suministrada al modelo;
y el arnés decide si muestrear, terminar o truncar la respuesta.
\begin{figure}[H]\centering
\begin{tikzpicture}[font=\small,>=Stealth,
box/.style={draw=azul,rounded corners,fill=azulclaro,align=center,text width=3cm,minimum height=1cm}]
\node[box] (p) at (0,0) {Prefijo\,(prompt)\\$x_{1:k}$};
\node[box] (q) at (4,0) {Distribución condicional\\$q_\theta(\cdot\mid h)$};
\node[box] (t) at (8,0) {Muestreo de $X_{k+1}$\\+ condición de parada};
\draw[->,thick] (p)--(q);
\draw[->,thick] (q)--(t);
\draw[->,thick] (t.south) -- ++(0,-.7) -| node[pos=.25,below]{si continúa: actualizar $h$} (p.south);
\end{tikzpicture}
\caption{El modelo calcula condicionales; el arnés transforma condicionales en un algoritmo de generación.}
\end{figure}
En particular, una cota de $n$ tokens impuesta externamente produce
una distribución de longitudes diferente de la ley generativa sin truncamiento.

\subsubsection*{Puente hacia la práctica en R y Python}

La estrategia pedagógica separa la representación de texto, la construcción
de una matriz de conteos y su interpretación probabilística. El módulo
\nolinkurl{modulo_codigo_s01.tex} proporciona un corpus ilustrativo, código base R
sin dependencias, validación de sumas por fila, cálculo de entropía cruzada
y una implementación didáctica de BPE con desempates deterministas.
Las arquitecturas recurrentes, convolucionales y \emph{Transformer} se
mencionan solo como extensiones: su tratamiento detallado corresponde
a sesiones posteriores.

% ===== Fin fuente =====

\subsection{Cierre del laboratorio matemático: predicciones que el código debe verificar}
\phantomsection\label{s01:lab:mat:cierre}
\begin{enumerate}
\item Sobre \texttt{abacabbacb}, la fila suavizada de $a$ es
$\left(\frac17,\frac37,\frac37\right)$, no $(1/6,\ldots)$:
$N_a=4$, $K=3$ y $\alpha=1$.
\item Para \texttt{abca}, la \emph{probabilidad condicionada en el primer
símbolo} es el producto de tres transiciones. La probabilidad \emph{conjunta}
incluye la masa inicial declarada. Mezclar ambos protocolos cambia la PPL.
\item La identidad $H(p,q)=H(p)+D_{\mathrm{KL}}(p\Vert q)$ es una igualdad
matemática (con convenciones de soporte), no una afirmación de que un modelo
entrenado estima correctamente $p$.
\item En una evaluación reproducible, la unidad de resampling es el documento;
no se deben remuestrear caracteres relacionados como observaciones independientes.
\end{enumerate}
\begin{puente}[Puente al laboratorio computacional]
La teoría define qué se calcula. En el laboratorio computacional de esta sesión
se declara el corpus, se fijan reglas de preprocesamiento y se generan
\nolinkurl{resultados_python.csv} y \texttt{resultados\_R.csv}; la concordancia
numérica de ambos es una prueba de implementación, no de generalización.
\end{puente}

% ===== Fin fuente =====


% ===== Fuente modulos/30_laboratorio_computacional.tex =====

\subsection{Laboratorio computacional: de las ecuaciones al reporte}
\phantomsection\label{s01:lab:comp}
\begin{idea}[Pregunta de trabajo]
\emph{¿Una tabla de transiciones estimada con lenguaje técnico asigna mayor
probabilidad a bitácoras técnicas reservadas que una tabla estimada con prosa general?}
Responder esta pregunta construida en pequeño ayuda a comprender qué medir,
pero no demuestra nada sobre un tokenizador moderno ni sobre la población real
de documentos industriales.
\end{idea}
\subsubsection*{L1. Inspeccionar el corpus antes de ajustar el modelo}
El archivo \nolinkurl{experimentos/corpus_didactico.csv} contiene ocho documentos
generales de entrenamiento, ocho técnicos de entrenamiento y seis técnicos
de prueba; todas las frases fueron \textbf{redactadas para esta clase}.
\textbf{No se descargaron informes de empresas ni se entrenó un LLM.}
Las dos condiciones usan el mismo vocabulario de caracteres establecido
\emph{solo con entrenamiento} y el mismo valor $\alpha=1$; los datos
de prueba nunca intervienen en el ajuste ni en la selección del modelo.
Se usan minúsculas y se quitan espacios de los extremos; el archivo se almacena
en UTF-8 con caracteres precompuestos, sin normalización adicional.

\subsubsection*{L2. Traducir ecuaciones a un algoritmo comprobable}
Para cada documento $z=z_1\cdots z_n$, se añade un estado inicial
\texttt{BOS} y se predice también un símbolo \texttt{EOS}. La probabilidad
predictiva de bigramas es
\[
\widehat P(v\mid u)=\frac{C(u,v)+\alpha}
{\sum_w C(u,w)+\alpha|V_{\mathrm{salida}}|},\qquad \alpha=1.
\]
La pérdida en bits para el documento completo es
\[
B(z)=-\sum_{t=1}^{n+1}\log_2 \widehat P(z_t\mid z_{t-1}),
\qquad \operatorname{BPC}(z)=\frac{B(z)}{n}.
\]
Aquí $z_0=\mathrm{BOS}$ y $z_{n+1}=\mathrm{EOS}$. El denominador
cuenta los puntos de código del texto original (\emph{no} EOS).
\textbf{Observación:} el alfabeto es el mismo y los modelos difieren en
los datos de entrenamiento, no en el algoritmo de tokenización.

\subsubsection*{L3. Ejecutar en dos lenguajes y localizar los resultados}
Desde la carpeta de la sesión S01:
\begin{verbatim}
python codigo/representacion.py
python codigo/s01.py
python codigo/experimento_s01.py
# En un equipo donde este instalado R:
Rscript codigo/representacion.R
Rscript codigo/s01.R
Rscript codigo/experimento_s01.R
\end{verbatim}
Los programas nuevos exportan los resultados a
\nolinkurl{experimentos/resultados_python.csv} y
\nolinkurl{experimentos/resultados_R.csv}. Python exporta también
\nolinkurl{experimentos/resumen_python.json}. Si falta \texttt{Rscript},
la ejecución R se deja \textbf{PENDIENTE}; no se informa una comparación
entre lenguajes como completada.

\subsubsection*{L4. Interpretar un resultado ejecutado, no extrapolarlo}
\begin{ejemplo}[Salida verificada por Python en este corpus sintético]
Se obtuvieron seis diferencias pareadas
$\Delta_i=\operatorname{BPC}_{\rm tecnico}(z_i)-
\operatorname{BPC}_{\rm general}(z_i)$; su promedio fue
$\bar\Delta\approx-0.54025$ bits por carácter.
Por enumeración exhaustiva de las $6^6=46656$ remuestras bootstrap
pareadas con reemplazo, el intervalo percentil nominal del 95\% fue
$[-0.66771,-0.42038]$.
El intervalo \emph{describe el comportamiento de un procedimiento de
remuestreo sobre seis frases construidas}, no la incertidumbre científica
de una población de reportes: no hay muestreo representativo, comparación
con modelos neuronales ni validación externa.
\end{ejemplo}
\subsubsection*{L5. De algoritmo a evidencia: registros obligatorios}
El reporte contiene la versión del programa, datos y licencia,
partición por documento, protocolo BOS/EOS, alfabeto, unidades,
valores por documento, método del intervalo, limitaciones y
una conclusión que no exceda lo medido. Cada cambio del corpus,
semillas, modelado o tokenización requiere regenerar todas las tablas.
El \hyperref[s01:lab:ic10]{microexperimento IC10} evalua, por separado,
la segmentación aprendida por dos tokenizadores BPE.
\ifdefined\labcompstandalone
La consigna del taller exige un reporte corto; el seminario de
investigación explica que falta
\else
La \hyperref[s01:mod:taller]{consigna del taller} exige un reporte corto;
el \hyperref[s01:mod:invest]{seminario de investigación} explica qué falta
\fi
para convertir una demostración en un estudio candidato a publicación.

% ===== Fin fuente =====


% ===== Fuente modulos/32_microexperimento_ic10.tex =====

\subsection{Microexperimento IC10: cambiar realmente el tokenizador}
\phantomsection
\label{s01:lab:ic10}
\begin{idea}[Ahora si cambiamos la representacion, no el estimador de probabilidades]
En el laboratorio anterior se compararon \emph{dos bigramas de caracteres}
con distinta procedencia de datos. Aquello ilustra estimacion de
probabilidades. Aqui comparamos dos \emph{tokenizadores BPE} que aprenden
fusiones diferentes, uno en prosa general y otro en prosa tecnica, con el
\textbf{mismo presupuesto de 20 fusiones}. Esta es una primera medicion de
la variable de representacion que estudia IC10, sin afirmar que un
modelo de lenguaje mejora por ello.
\end{idea}
\subsubsection*{Reglas transparentes}
La unidad inicial de cada palabra es un punto de código Unicode; se añade
un marcador \texttt{</w>} para no fusionar entre palabras. Cada paso combina
el par de unidades adyacentes mas frecuente del corpus de entrenamiento
(con desempate lexicografico). Las reglas se aplican, en el orden en que
se aprendieron, a las \textbf{mismas seis frases tecnicas de prueba}.
La separacion entrenamiento/prueba es por documento. Se conserva la
normalizacion explicita: minusculas, recorte externo de espacios y
separacion por blancos. No se eliminan acentos, ni se usa el texto de
prueba para elegir fusiones. Este BPE didactico no es el tokenizador de
un modelo neuronal preentrenado.
\subsubsection*{Resultado verificado en Python, datos sintéticos}
\begin{center}\begin{tabular}{lrrr}\toprule
Tokenizador & Palabras de prueba & Tokens en prueba & Tokens/palabra\\\midrule
BPE entrenado con prosa general & 35 & 148 & 4.228571\\
BPE entrenado con texto tecnico & 35 & 127 & 3.628571\\\bottomrule
\end{tabular}\end{center}
La diferencia agregada es $127-148=-21$ tokens y
$3.628571-4.228571=-0.600000$ tokens por palabra, bajo estas reglas.
La tabla por documento se genera en
\nolinkurl{experimentos/resultados_ic10_tokenizacion.csv}; las veinte
fusiones de cada modelo, en \nolinkurl{experimentos/fusiones_ic10.csv}.
\begin{verificacion}[Interpretacion responsable]
El resultado no permite afirmar que un modelo general tenga peor
perplejidad ni que el tokenizador especializado mejore una tarea aplicada.
Para contrastar \textbf{bits por caracter} entre tokenizadores se
necesita asignar probabilidades normalizadas y compatibles a las mismas
secuencias (incluyendo convenciones de fronteras y fin), y evaluar en
textos reservados. Menos tokens significa \emph{otra segmentación}; no
implica por si solo menor entropia cruzada ni mejor generalizacion.
\end{verificacion}
\subsubsection*{Ejecución y reporte}
\begin{verbatim}
python codigo/ic10_tokenizadores.py
# En un equipo con R instalado:
Rscript codigo/ic10_tokenizadores.R
python verificaciones/comparar_r_python.py
\end{verbatim}
El estudiante reporta las dos comparaciones por separado:
\textbf{A)} corpus de entrenamiento frente a BPC de bigramas por
caracteres; \textbf{B)} entrenamiento del tokenizador frente a tokens por
palabra. Ambas son ilustrativas y no equivalen a un estudio industrial.

% ===== Fin fuente =====


% ===== Fuente modulos/35_codigo_detallado.tex =====

\subsection{Protocolos complementarios de código y BPE}
\phantomsection

% ===== Fuente fuentes/modulo_representacion.tex =====
\subsection{Primeros pasos: tokenizar, formar ventanas y codificar}
\subsubsection{El ejemplo de la guía en código sencillo}
El siguiente ejemplo es independiente de Keras o de cualquier modelo neuronal. Solo pretende que el estudiante vea la misma operación en Python y R sobre la frase \texttt{el gato mira al perro}. La unidad es la \emph{palabra separada por espacios}; no se presenta como un tokenizador moderno completo.
\paragraph{Python (biblioteca estándar)}
\begin{lstlisting}[language=Python]
texto = "el gato mira al perro"
tokens = texto.split()
bigramas = list(zip(tokens, tokens[1:]))
vocabulario = list(dict.fromkeys(tokens))
one_hot_gato = [int(w == "gato") for w in vocabulario]
assert len(tokens) == 5 and len(bigramas) == 4
assert sum(one_hot_gato) == 1
\end{lstlisting}
\paragraph{R (sin paquetes externos)}
\begin{lstlisting}[language=R]
texto <- "el gato mira al perro"
tokens <- strsplit(texto, ' ', fixed=TRUE)[[1]]
bigramas <- lapply(seq_len(length(tokens)-1L),
                   function(i) tokens[i:(i+1L)])
vocabulario <- unique(tokens)
one_hot_gato <- as.integer(vocabulario == 'gato')
stopifnot(length(tokens)==5L, length(bigramas)==4L,
          sum(one_hot_gato)==1L)
\end{lstlisting}
\noindent Compare los resultados, no únicamente la apariencia del código. Los programas completos \nolinkurl{codigo/representacion.py} y \nolinkurl{codigo/representacion.R} muestran además trigramas, bolsa de palabras e identificadores.
\subsubsection{Lo que la representación no demuestra}
Un vector one-hot identifica un token, pero no expresa semejanza semántica ni es una probabilidad predicha. Un embedding es un vector aprendido que estudiaremos en S02. Una bolsa de palabras conserva frecuencias, no orden; y un vector por documento puede corresponder a una tarea de clasificación distinta de la predicción autorregresiva.
\subsubsection{Cómo ejecutar y comparar}
\begin{lstlisting}
python codigo/representacion.py
python codigo/s01.py
python verificaciones/verificar_s01.py
Rscript codigo/representacion.R
Rscript codigo/s01.R
\end{lstlisting}
\noindent Los comandos R requieren \texttt{Rscript} en el sistema. Si no está disponible, las instrucciones se entregan con una prueba pendiente, y nunca se atribuye ejecución a R sin observarla.

% ===== Fin fuente =====


% ===== Fin fuente =====


% ===== Fuente modulos/codigo_detalle_heredado.tex =====
\subsection{Datos y convenciones compartidas}
La cadena de entrenamiento es \texttt{abacabbacb}; el conjunto de prueba
\texttt{abca}; el alfabeto $V=\{a,b,c\}$. La matriz de bigramas usa
filas para el símbolo precedente y columnas para el siguiente:
\[
C=\begin{pmatrix}0&2&2\\2&1&0\\1&1&0\end{pmatrix}.
\]
Se adoptan logaritmos base dos para los resultados en bits y la
probabilidad inicial $p(X_1=i)=1/3$ para la evaluación conjunta del
microejemplo. Los conteos se forman solo con entrenamiento; se evalúa
prueba sin alterarlos.
\subsubsection{R: matriz y suavizado}
\begin{lstlisting}[language=R]
train <- strsplit("abacabbacb", "", fixed=TRUE)[[1]]
V <- c("a", "b", "c")
C <- matrix(0L, 3, 3, dimnames=list(V,V))
for (i in seq_len(length(train)-1L))
    C[train[i],train[i+1L]] <- C[train[i],train[i+1L]]+1L
P <- sweep(C+1, 1, rowSums(C)+3, "/")
stopifnot(all(abs(rowSums(P)-1)<1e-12))
\end{lstlisting}
\subsubsection{Python: matriz y suavizado}
\begin{lstlisting}[language=Python]
train = "abacabbacb"
V = "abc"
C = {a: {b: 0 for b in V} for a in V}
for a, b in zip(train, train[1:]):
    C[a][b] += 1
P = {a: {b: (C[a][b] + 1)/(sum(C[a].values()) + 3)
          for b in V} for a in V}
assert all(abs(sum(row.values())-1) < 1e-12
           for row in P.values())
\end{lstlisting}
Ambos programas deben obtener
$P(a\to b)=3/7$, $P(b\to c)=1/6$ y $P(c\to a)=2/5$.
Una revisión útil consiste en contrastar no solo números redondeados,
sino también suma de conteos, orden del vocabulario y convenciones
sobre contextos sin observaciones.
\subsection{De condicionales a pérdida, PPL y BPC}
La probabilidad conjunta de la prueba es
\[
q(\texttt{abca})=\frac13\frac37\frac16\frac25=\frac1{105}.
\]
\begin{lstlisting}[language=Python]
from math import log2
bits = log2(105)
bpc = bits / 4
ppl = 2 ** bpc
assert abs(bits - 6.714245517666122) < 1e-10
\end{lstlisting}
\begin{lstlisting}[language=R]
bits <- log2(105)
bpc <- bits / 4
ppl <- 2^bpc
stopifnot(abs(bits - 6.714245517666122) < 1e-10)
\end{lstlisting}
Si se omite $\pi(a)=1/3$ en la verosimilitud y se promedian las
cuatro posiciones de la cadena, se está mezclando una pérdida
\emph{condicional} con la tasa conjunta. Es preferible reportar
explícitamente la convención de evaluación.
\subsection{BPE con desempate definido}
Se procesan palabras con frecuencias
\emph{casa} (5), \emph{casas} (2), \emph{cosa} (6),
\emph{cosas} (3), \emph{caso} (4).
El marcador \texttt{<eow>} se incluye al final de cada palabra.
El conteo de pares se pondera por frecuencia del tipo léxico;
los empates se resuelven por orden Unicode ascendente sobre el par
reconstruido, y los reemplazos se ejecutan de izquierda a derecha
sin solapamiento.
\begin{center}\small
\begin{tabular}{clrr}\toprule Paso & Par & Reemplazos & Longitud\\\midrule
0&---&---&105\\1&$a+s$&16&89\\2&$a+\mathrm{fin}$&11&78\\
3&$c+as$&11&67\\4&$c+o$&9&58\\5&$co+s$&9&49\\
6&$cos+a\mathrm{fin}$&6&43\\\bottomrule
\end{tabular}
\end{center}
Una fila del registro tiene dos contadores posibles: frecuencia bruta
del par y número de sustituciones efectivas. Cuando existen
solapamientos no tienen por qué coincidir. La comprobación correcta es
\[
\text{longitud anterior}-\text{longitud nueva}
=\text{número de sustituciones efectivas ponderadas}.
\]
El código completo (sin omisiones) está en
\texttt{codigo/s01.py} y \texttt{codigo/s01.R}. Se evitan aquí bloques
extensos de código repetitivo para mantener el documento enfocado en
sus invariantes y en las relaciones matemáticas.
\subsection{Extensión a investigación: cómo empezar sin contaminar la prueba}
\subsubsection{Un protocolo de documentos reales}
\begin{enumerate}
\item Seleccionar una fuente general y otra técnica con licencia declarada,
sin suponer que el catálogo del documento de investigación contiene
archivos de corpus ya descargados.
\item Congelar particiones por documento o grupo de documentos relacionados
antes de aprender el vocabulario BPE.
\item Entrenar tokenización y $n$-gramas solo con entrenamiento;
elegir suavizado y tamaños de vocabulario con validación.
\item Registrar tasa de tokens por palabra, longitud, bits por carácter
y, cuando corresponda, pérdida por token compatible.
\item Presentar incertidumbre por documento o grupo, no por supuesto de
independencia entre todos los tokens de un mismo documento.
\end{enumerate}
\subsubsection{Separación entre ejemplo y evidencia}
La matriz, el ejemplo de KL, la tabla BPE y el contraste entre PPL
son \textbf{ILUSTRATIVOS con aritmética verificada}. Una mejora en
reportes técnicos se considerará \textbf{REAL} solo después de una
corrida documentada sobre datos admisibles, con código y control de
fuga. Ninguna de las cifras de este módulo sustituye el futuro
protocolo de investigación.
\subsubsection{Pruebas mínimas automatizadas}
\begin{enumerate}
\item La suma de las transiciones es igual al número de símbolos
menos uno.
\item Cada fila de probabilidades de una matriz válida suma uno.
\item La prueba suavizada tiene probabilidad conjunta $1/105$.
\item La divergencia KL del ejemplo vale $1/4$ bit.
\item El procedimiento BPE pasa de longitud 105 a 43 en seis pasos.
\end{enumerate}
Estas condiciones están programadas en
\nolinkurl{verificaciones/verificar_s01.py}, que usa las mismas
funciones públicas del módulo Python.

% ===== Fin fuente =====


% ===== Fuente modulos/50_taller_contenido.tex =====

\subsection{Taller autónomo con reporte técnico breve}
\phantomsection\label{s01:mod:taller}
\begin{taller}[Finalidad y entrega]
Demostrar, con cálculos y archivos regenerables, cómo las decisiones de
representación afectan a un pequeño modelo de lenguaje. Entregar un
\textbf{reporte de 2--3 páginas} (objetivo, datos, métodos, resultados,
incertidumbre, limitaciones y relación con IC10), más scripts y CSV.
\textbf{Todo es trabajo autónomo}; no se presuponen las horas del laboratorio
dentro de la clase presencial. En S01 se conservan como hito de investigación
\textbf{dos áreas de interés}, cada una con un párrafo.
\end{taller}
\subsubsection*{Actividad 1 [Representación]: una frase, cuatro objetos}
Segmente ``el gato mira al perro'' en palabras y caracteres; construya
unigramas, bigramas y trigramas; contraste bolsa de palabras con secuencia y
escriba el vector one-hot del token \texttt{gato}. Repita el primer paso con
una frase técnica. \emph{Autoverificación}: cinco palabras, cuatro pares,
tres trigramas en la primera oración.
\subsubsection*{Actividad 2 [Cálculo]: del conteo a la predicción}
Con \texttt{abacabbacb} construya la matriz de transiciones, derive el
estimador MLE e introduzca $\alpha=1$. Para \texttt{abca}, calcule la
probabilidad de transición y la conjunta con $p(X_1=a)=1/3$.
\emph{Autoverificación}: $P(b\mid a)=3/7$; probabilidad conjunta $1/105$.
\subsubsection*{Actividad 3 [Demostración]: medir lo que se predice}
Demuestre la identidad entre entropía cruzada y KL con las condiciones de
soporte; calcule KL de $(1/2,1/4,1/4)$ respecto de
$(1/4,1/2,1/4)$. Para las probabilidades observadas
$(1/2,1/4,1/8,1/2)$, calcule PPL. Compare dos tokenizadores con
$(T_A,\mathrm{PPL}_A)=(100,3.2)$ y $(T_B,\mathrm{PPL}_B)=(30,25)$
sobre un mismo texto de 100 caracteres, indicando qué supuestos permiten
interpretar el BPC. \emph{Autoverificación}: KL $=0.25$ bits y
PPL $=2^{7/4}$.
\subsubsection*{Actividad 4 [Algoritmos]: BPE, R y Python}
Ejecute seis fusiones BPE sobre la tabla \emph{casa, casas, cosa, cosas,
caso} con las frecuencias y reglas indicadas en la guía. Ejecute después:
\begin{verbatim}
python codigo/s01.py
python codigo/experimento_s01.py
python codigo/ic10_tokenizadores.py
Rscript codigo/s01.R
Rscript codigo/experimento_s01.R
Rscript codigo/ic10_tokenizadores.R
\end{verbatim}
Compruebe las sumas de las filas, las longitudes BPE y las seis filas
pareadas del CSV. Compare también, sobre las mismas seis frases
de prueba, el BPE general con el tecnico mediante 20 fusiones en ambos;
produzca \nolinkurl{resultados_ic10_tokenizacion.csv} y explique las diferencias.
Si no tiene R, indique \textbf{PENDIENTE} en el reporte,
sin inventar que los números se reprodujeron en ambos lenguajes.
\subsubsection*{Actividad 5 [Aplicación e investigación]: entregar evidencia}
Use \nolinkurl{experimentos/corpus_didactico.csv} para reconstruir
\nolinkurl{experimentos/resultados_python.csv}; complete la
\textbf{plantilla de reporte} en \nolinkurl{reportes/reporte_s01.tex}.
Incluya la tabla BPE con tokens por palabra y distinga la comparación
de tokenizadores de la comparación de predictores por caracteres.
Interprete la diferencia media de BPC y el intervalo bootstrap por
\emph{documento}, explicando por qué ni el tamaño de muestra ni el origen
del corpus permiten inferencia para textos industriales reales.
Proponga una versión de IC10 que sí sea falsable con documentos y
comparaciones externas. Finalmente, formule dos áreas de interés,
con un párrafo justificado por cada una.
\begin{verificacion}[Lista de control antes de entregar]
\begin{enumerate}
\item Los datos no se mezclan entre ajuste y prueba; se identifica cada
unidad de observación.
\item Se reportan todas las probabilidades de la cadena, incluido EOS,
y se cuenta explícitamente la unidad del denominador BPC.
\item Cada resultado se etiqueta como \emph{calculado}, \emph{ejecutado
con datos sintéticos} o \emph{pendiente de validación empírica}.
\item Se entregan código, CSV, versión de lenguaje, interpretación,
limitaciones y dos áreas de interés.
\end{enumerate}
\end{verificacion}
\subsubsection*{Extensiones opcionales, fuera del núcleo S01}
Calcule intervalo de confianza mediante bootstrap jerárquico con varios
documentos por fuente; incorpore una evaluación BPC de modelos probabilisticos sobre tokens
BPE con distribuciones comparables; compare
vocabularios fijados por entrenamiento con distintos tamaños; pruebe
suavizado $\alpha\in\{1/2,1,2\}$ exclusivamente en validación; diseñe una
réplica externa para dos dominios técnicos con control de licencias.

% ===== Fin fuente =====


% ===== Fuente modulos/40_investigacion_y_q1.tex =====

\subsection{De la utilidad práctica al objetivo de publicación Q1}
\phantomsection\label{s01:mod:invest}
\begin{investigacion}[Una meta ambiciosa con etapas verificables]
La trayectoria de los tres cursos tendrá como objetivo
\textbf{preparar un manuscrito competitivo para una revista del primer
cuartil (Q1) de una categoría pertinente}. Q1 es una clasificación
bibliométrica dependiente de \emph{base, categoría y año}; no es garantía
de aceptación ni un sustituto de la calidad del método.
En \textbf{este primer curso} se exige un protocolo, una línea base,
resultados preliminares honestos y un informe técnico reproducible.
No se exigirá enviar el artículo en S12. La novedad, el alcance externo y
las validaciones adicionales corresponden a la trayectoria de investigación
posterior y a la selección de revista.
\end{investigacion}
\subsubsection*{Qué problema puede llegar a interesar a una revista exigente}
Los experimentos de \hyperref[s01:lab:comp]{bigramas de caracteres}
y \hyperref[s01:lab:ic10]{tokenización BPE} ensenan preguntas distintas:
el primero compara perdida predictiva con tokenización fija; el segundo
mide longitud tokenizada, sin evaluar probabilidades de un modelo.
Ninguno constituye novedad publicable. Una ruta posible desde IC10 es estudiar
\emph{bajo qué condiciones} la segmentación de textos técnicos en español
aumenta el costo de representación y, cuando se controla el modelo de
probabilidad y su capacidad, altera la calidad predictiva o el desempeño
de tareas del dominio. Es necesario contrastar la hipótesis con la literatura
vigente y establecer qué mecanismo o efecto no ha sido caracterizado.
\subsubsection*{Qué garantías faltarán para un artículo candidato a Q1}
\begin{enumerate}
\item \textbf{Vacío verificable en la literatura:} revisión sistemática de trabajos
pertinentes; plantear una contribución diferenciada de la ingeniería rutinaria.
\item \textbf{Datos pertinentes y auditables:} varias fuentes y dominios si la
pregunta lo exige, permisos y licencias, documentos fechados y prueba reservada.
\item \textbf{Comparaciones fuertes:} tokenizadores con presupuesto de vocabulario
controlado y modelos probabilísticos bien definidos; baselines de carácter,
subpalabra y modelo preentrenado; coste/latencia donde corresponda.
\item \textbf{Diseño reproducible:} normalización definida, corpus separados por
documento, evidencia de ausencia de duplicados/fuga, experimentos pareados,
varias semillas y control de selección múltiple.
\item \textbf{Validez externa:} réplica en otros documentos o instituciones
cuando los permisos lo permitan; ablationes, análisis de errores y discusión
de los resultados nulos o negativos.
\item \textbf{Comunicación científica:} figuras regeneradas desde los datos,
reporte de tamaño de efecto e incertidumbre, materiales compartibles y
justificación del alcance real de la contribución.
\end{enumerate}
\begin{puente}[De la S01 a las sesiones posteriores]
\textbf{S01:} explorar dos áreas y practicar el reporte. \textbf{S02--S04:}
reducir candidatos y comprender mecanismos. \textbf{S05:} seleccionar el
problema con una rúbrica. \textbf{S06:} medir la línea base. \textbf{S07:}
prerregistrar la prueba. \textbf{S08--S11:} ejecutar, medir, analizar
sensibilidad y errores. \textbf{S12:} informe técnico, repositorio y revisión
por pares. \textbf{Cursos II--III:} experimentos completos, validación
externa y manuscrito sometible, con objetivo Q1 condicionado a la
aportación y al ajuste con la revista.
\end{puente}
\subsubsection*{Producto de S01}
\textbf{A entregar:} un párrafo por cada una de dos áreas de interés
(problemática, relevancia, datos candidatos y riesgo); no es aún la ficha
del problema definitivo. El pequeño informe del taller muestra que el
estudiante puede \emph{producir y criticar evidencia}; no constituye el
artículo científico.

% ===== Fuente fuentes/seminario_investigacion.tex =====
\subsection{Seminario de investigación: las preguntas comienzan hoy}

\begin{investigacion}[Meta intelectual de la última hora]
No se pretende seleccionar un tema definitivo ni entregar un experimento acabado.
Se pretende aprender a distinguir una demostración tecnológica de una pregunta
contrastable y a identificar \textbf{dos áreas de interés}, con un párrafo razonado
por área. La ficha formal será en S05; la línea base medida, en S06; y el
prerregistro, en S07. Esta pauta procede del documento de investigación del curso,
\emph{Integración e investigación aplicada}, capítulo 3 (2026).
\end{investigacion}

\subsubsection*{Problema inicial: tokenización de lenguaje técnico en español}

Una herramienta puede fragmentar \emph{empacadura}, \emph{revestidor},
\emph{pérdida de circulación} o \emph{9 5/8 pulgadas} de maneras distintas.
Esta observación sugiere una pregunta, pero aún no constituye una hipótesis
científica: ¿una tokenización especializada necesita menos unidades por término,
y qué costo tiene ese ahorro en una predicción probabilística razonable?

La matriz institucional identifica este hilo como \textbf{IC10}: penalización de
la tokenización en textos técnicos en español. Los observables propuestos son
\emph{tokens por palabra} y \emph{bits por carácter}; el catálogo vincula la
primera exploración con D01 y D02. El documento proporciona \emph{referencias de
fuentes de datos}, no un corpus descargado y preparado para la clase.

\begin{figure}[H]\centering
\begin{tikzpicture}[>=Stealth,font=\small,
 q/.style={draw=azul,fill=azulclaro,rounded corners,align=center,text width=3.2cm,minimum height=13mm}]
\node[q] (p) at (0,0) {\textbf{Pregunta}\\¿Qué se fragmenta?};
\node[q] (h) at (4.2,0) {\textbf{Hipótesis}\\un dominio cambia la tasa};
\node[q] (m) at (8.4,0) {\textbf{Medida}\\tokens/palabra; BPC};
\node[q] (d) at (8.4,-2.0) {\textbf{Datos}\\corpus con licencia};
\node[q] (v) at (4.2,-2.0) {\textbf{Validación}\\documentos reservados};
\node[q] (r) at (0,-2.0) {\textbf{Resultado}\\efecto e incertidumbre};
\draw[->,thick] (p)--(h);\draw[->,thick] (h)--(m);\draw[->,thick] (m)--(d);\draw[->,thick] (d)--(v);\draw[->,thick] (v)--(r);
\end{tikzpicture}
\caption{Estructura mínima de una pregunta experimental para la trayectoria de investigación.}
\end{figure}

\subsubsection*{Cuatro distinciones que hacen investigable la pregunta}

\textbf{1. Variable observable.} La fracción de palabras partidas en subunidades
no es idéntica a los bits por carácter. La primera caracteriza la representación;
la segunda resume una pérdida probabilística si el modelo permite computarla.
\textbf{2. Tratamiento.} Entrenar un nuevo BPE modifica la segmentación; cambiar
simultáneamente arquitectura, número de parámetros y corpus no permite atribuir
una mejora solo a la segmentación.
\textbf{3. Unidad experimental.} Debe ser, cuando sea posible, el documento o
el informe, no cada subfragmento del mismo archivo; de lo contrario,
los intervalos basados en tokens individuales sobrestiman la precisión.
\textbf{4. Generalización.} Un vocabulario reducido para diez términos elegidos a mano
no es evidencia de efecto promedio en un dominio. Es una demostración inicial.

\begin{ejemplo}[Cómo pasar de observación a hipótesis falsable]
\textbf{Observación ilustrativa:} un tokenizador general parte el término
\emph{revestidor} en más fragmentos que otro entrenado en textos de petróleo.
\textbf{Hipótesis operacional futura:} sobre informes técnicos reservados,
con normalización y tamaño de vocabulario previamente fijados,
la diferencia media pareada de tokens por palabra es mayor que cero.
\textbf{Evaluación:} informar media, intervalo de confianza por remuestreo de
documentos y posibles fallas en términos que no estaban en el vocabulario.
No se informa un resultado numérico porque los datos no han sido descargados
ni el experimento ejecutado.
\end{ejemplo}

\subsubsection*{Fuentes candidatas, derechos y riesgos}

\begin{center}\small
\begin{tabularx}{\linewidth}{@{}p{1.2cm}p{3.3cm}Xp{2.6cm}@{}}\toprule
ID & Fuente del catálogo & Uso inicial posible & Precaución\\\midrule
D01 & Wikipedia (eswiki) & Corpus amplio de artículos en español y temas científicos & Atribución y condiciones de reutilización según volcado/licencia\\
D02 & Proyecto Gutenberg & Obras en español para contraste con prosa técnica & Derecho de autor según país y obra; condiciones del portal\\
D03 & Volve, Equinor & Candidato posterior para texto y series de petróleo & Verificar presencia de texto, autorización y licencia de los archivos\\\bottomrule
\end{tabularx}
\end{center}
\textbf{Ubicaciones de origen:} \url{https://dumps.wikimedia.org/eswiki/},
\url{https://www.gutenberg.org/} y
\url{https://www.equinor.com/energy/volve-data-sharing}.
La fuente de la licencia debe verificarse antes de redistribuir; la existencia
de un catálogo no equivale a disponer de permiso para todo uso. La página de
condiciones de Gutenberg señala explícitamente la necesidad de examinar los
derechos en el país del usuario, además de los términos de la obra.

\subsubsection*{Conexión con Series de Tiempo y con otras áreas}

El objeto probabilístico $p(x_{1:T})$ aparece en ambos cursos. Una cadena de
Markov utiliza una matriz de transición; un modelo de lenguaje tabular cuenta
símbolos y estima esas transiciones. La similitud no implica que un LLM sea una
cadena homogénea finita: sus condicionales pueden depender de una ventana amplia
y de una función parametrizada por una red.

Además de IC10, la matriz incluye problemas que conectan texto y pronóstico:
PG02 (eventos de perforación), FI01 (tono de noticias y volatilidad) y EE01
(demanda eléctrica con eventos textuales). Son \emph{posibilidades de exploración},
no compromisos de ejecución en la S01. En trabajos de pronóstico con modelos
preentrenados habrá que controlar también la fuga temporal y lo que el modelo
pueda haber memorizado; una partición cronológica de nuestros archivos no elimina
por sí sola ese riesgo.

\subsubsection*{Guion de discusión de la última hora}

\begin{enumerate}
\item ¿Qué hipótesis distingue una comparación de tokenizadores de una
 comparación de modelos completos? Dar una variable que se pueda medir.
\item ¿Qué evidencia refutaría la hipótesis en vez de confirmarla?
\item ¿Qué datos hay hoy, cuáles requieren autorización y qué licencias
 pueden condicionar la investigación?
\item ¿Qué cambia si la unidad de análisis es un documento y no un token?
\item ¿Cómo se formularía una pregunta que conecte texto fechado con una
 serie temporal, evitando que la información futura entre en el entrenamiento?
\end{enumerate}
\begin{investigacion}[Hito autónomo S01]
Cada estudiante entrega \textbf{dos áreas de interés}, un párrafo argumentado
por área. En cada párrafo se solicita: relevancia del problema, una posible
fuente de datos, una variable que eventualmente podría medirse y una cautela.
No se exige entrenar un modelo ni comprometer una hipótesis definitiva.
Esta secuencia anticipa los cinco candidatos de S02 y la decisión formal de S05.
\end{investigacion}

\subsubsection*{Del curso a una evidencia reproducible}

El repositorio del proyecto conservará configuración de entornos, fuentes,
transformaciones, semillas, pruebas y salidas regenerables. Las tablas de resultados
no deben editarse a mano. En ejercicios con IA generativa deberá declararse su
uso; toda derivación matemática deberá ser defendible y toda cita verificable.
Los proyectos con datos personales requieren procedimientos de anonimización y,
si corresponde, autorizaciones. Una herramienta que responde correctamente a una
pregunta de demostración no constituye una evidencia científica.

% ===== Fin fuente =====


% ===== Fin fuente =====


% ===== Fuente modulos/60_suplemento_contenido.tex =====
\phantomsection\label{s01:mod:sup}
\subsection{Distribuciones finitas y secuencias}
\subsubsection{El símplex y el soporte}
Sea $V=\{1,\ldots,K\}$. El símplex es
$\Delta_{K-1}=\{p\in\mathbb R^K:p_i\ge0,\sum_i p_i=1\}$.
Su interior relativo exige $p_i>0$. El soporte es
$\operatorname{supp}(p)=\{i:p_i>0\}$. La convención
$0\log0=0$ está motivada por $\lim_{x\downarrow0}x\log x=0$.
Cuando $q_i=0$ y $p_i>0$, $-p_i\log q_i=+\infty$.
\subsubsection{Regla de la cadena con cuidado del soporte}
\textbf{Proposición.} Si $p(x_{1:T})>0$, entonces
\[p(x_{1:T})=p(x_1)\prod_{t=2}^T p(x_t\mid x_{1:t-1}).\]
\textit{Demostración.} Para todo prefijo del soporte, la definición de
condicional permite escribir
$p(x_{1:t})=p(x_{1:t-1})p(x_t\mid x_{1:t-1})$.
Al iterar, los factores de prefijos se cancelan telescópicamente.\hfill$\square$
Para prefijos de probabilidad cero, las condicionales pueden elegirse por
convención; la representación de la ley conjunta sigue siendo válida donde
está definida. Esto evita atribuir contenido probabilístico a expresiones
condicionales que no son identificables por datos.
\subsubsection{Markov no es una consecuencia de la factorización}
La hipótesis de orden $m$ afirma
$X_t\perp X_{1:t-m-1}\mid X_{t-m:t-1}$ bajo las condiciones de la ley.
La notación de independencia condicional es una propiedad de la distribución;
una ventana de longitud $m$ en un algoritmo es una \emph{decisión de modelado}.
Por ejemplo, si $X_1,X_2$ son monedas justas independientes y $X_3=X_1$,
el modelo de orden uno pierde información relevante para el último paso.
\subsubsection{Número de parámetros de una tabla condicional}
Para $K$ símbolos y una historia de longitud $m$ (por tanto, $K^m$ contextos posibles), cada contexto posee $K-1$
grado de libertad; la tabla completa usa $K^m(K-1)$ parámetros libres.
Este conteo presupone que cada contexto se trata independientemente, que no
hay simetrías ni parámetros compartidos, y que todas las filas pueden variar.
Las redes neuronales comparten parámetros y no se cuentan de esta forma.

\subsection{Estimación, concavidad y modelos predictivos}
\subsubsection{Existencia del estimador de máxima verosimilitud}
Sea $c_i\ge0$, $N=\sum c_i>0$ y
$\ell(p)=\sum_{i:c_i>0}c_i\ln p_i$. Es una función cóncava en el
símplex: el Hessiano en coordenadas positivas observadas es diagonal con
entradas $-c_i/p_i^2$. Para maximizar, la desigualdad de Gibbs o el
multiplicador de Lagrange en el soporte observado dan
$\widehat p_i=c_i/N$. A los índices con $c_i=0$ se les asigna masa nula.
Si se impusiera artificialmente $p_i>0$ para todos, el supremo podría no alcanzarse.
\subsubsection{Dirichlet y posterior: derivación completa}
Para $\alpha_i>0$ la densidad de $p\sim\mathrm{Dirichlet}(\alpha)$ es
\[
f(p)=\frac{\Gamma(A)}{\prod_i\Gamma(\alpha_i)}
\prod_{i=1}^Kp_i^{\alpha_i-1},\qquad A=\sum_i\alpha_i.
\]
Con conteos $c$, la verosimilitud es proporcional a $\prod p_i^{c_i}$.
Por Bayes la posterior es Dirichlet$(\alpha_1+c_1,\ldots,\alpha_K+c_K)$.
La identidad de momentos de Dirichlet es
$\mathbb E[p_i\mid c]=(\alpha_i+c_i)/(A+N)$.
Se demuestra integrando, pues la integral de $p_i\prod_j p_j^{\beta_j-1}$
con la constante de normalización es $\beta_i/(\sum_j\beta_j)$.
Así, la predicción posterior de un símbolo nuevo es la media posterior;
no coincide necesariamente con la moda posterior.
\subsubsection{Suavizado como estimación de compromiso}
Cuando $\alpha_i=\alpha$,
\[
\frac{c_i+\alpha}{N+K\alpha}
=\frac{N}{N+K\alpha}\frac{c_i}{N}+
\frac{K\alpha}{N+K\alpha}\frac1K.
\]
Esto muestra una retracción hacia la distribución uniforme, no una garantía
de buena generalización. Si $K$ crece, la masa agregada aumenta.
\subsubsection{Selección y evaluación separadas}
La elección de $\alpha$, $n$ o las fusiones BPE sobre prueba destruye la
interpretación de prueba como evidencia final independiente. La recomendación
es separar por documento antes de ajustar cualquier representación y realizar
validación de hiperparámetros solo sobre el conjunto de validación.

\subsection{Entropía y divergencia KL}
\subsubsection{Entropía como esperanza}
Para una distribución finita $p$, la información de $i$ es
$I_p(i)=-\log_2p_i$ y $H_2(p)=\mathbb E_p[I_p(X)]$.
Con base natural se mide en nats. La relación es
$H_2(p)=H_e(p)/\ln2$. La entropía mide incertidumbre bajo $p$;
la entropía cruzada $H(p,q)$ mide pérdida media al codificar datos de $p$
con probabilidades $q$.
\subsubsection{Límites y casos de igualdad}
\textbf{Proposición.} Para $K$ categorías,
$0\le H_2(p)\le\log_2K$. La primera desigualdad se obtiene porque
$-\log_2p_i\ge0$ cuando $0<p_i\le1$; la igualdad ocurre para masas
concentradas en un solo evento. La segunda se deduce de
$D_{\rm KL}(p\|u)=\log_2K-H_2(p)\ge0$, donde $u$ es uniforme.
La igualdad ocurre exactamente para $p=u$.
\subsubsection{Desigualdad de Gibbs con masas nulas}
\textbf{Teorema.} $D_{\rm KL}(p\|q)\ge0$, con igualdad si y solo si $p=q$.
\textit{Demostración.} Si existe $i$ con $p_i>0=q_i$, la divergencia es
infinita. En el otro caso, usando $\ln z\le z-1$,
\[
-D_{\rm KL}(p\|q)=\sum_{i:p_i>0}p_i\ln\frac{q_i}{p_i}
\le\sum_{i:p_i>0}(q_i-p_i)\le0.
\]
Para la igualdad, la desigualdad logarítmica exige $q_i=p_i$ en todo el
soporte y la última desigualdad exige que $q$ no tenga masa fuera de él.
Concluimos $p=q$.\hfill$\square$
\subsubsection{Entropía cruzada y KL}
Si los soportes son compatibles,
\[
H(p,q)=-\sum_ip_i\log q_i
=-\sum_ip_i\log p_i+\sum_ip_i\log\frac{p_i}{q_i}
=H(p)+D_{\rm KL}(p\|q).
\]
Por tanto $q=p$ minimiza la entropía cruzada sobre \emph{todas} las
distribuciones posibles. Si se optimiza una familia restringida $q_\theta$,
se obtiene el mejor elemento de esa familia según el criterio, no
necesariamente la distribución real.
\subsubsection{KL no es una métrica}
Para $p=(3/4,1/4)$ y $q=(1/2,1/2)$,
\[
D_{\rm KL}(p\|q)=1-H_2(p)\approx0.188722\text{ bits},
\quad D_{\rm KL}(q\|p)=\tfrac12\log_2(4/3)\approx0.207519.
\]
La falta de simetría basta para concluir que KL no es una métrica.
Es frecuente usar una divergencia simétrica como Jensen--Shannon:
$\mathrm{JS}(p,q)=\tfrac12D_{\rm KL}(p\|m)+\tfrac12D_{\rm KL}(q\|m)$,
con $m=(p+q)/2$. No se necesita JS en la S01, pero permite situar KL
entre otros objetos de la teoría de la información.
\subsubsection{KL entre distribuciones sobre secuencias}
Para $p,q$ sobre $X_{1:T}$, con soportes compatibles,
\[
D_{\rm KL}(p(x_{1:T})\|q(x_{1:T}))=
\sum_{t=1}^T\mathbb E_p\left[
D_{\rm KL}\bigl(p(X_t\mid X_{<t})\|q(X_t\mid X_{<t})\bigr)\right].
\]
Se deriva factorizando el cociente de densidades o masas mediante la regla
de la cadena y conmutando la suma finita y la esperanza. La fórmula es una
versión de la regla de la cadena para KL, y enlaza la modelación de
secuencias con las diferencias de pérdida condicional.

\subsection{Perplejidad, codificación y tokenización}
\subsubsection{La perplejidad como transformación de la pérdida}
Para una realización observada y $r_t=q(x_t\mid x_{<t})>0$,
\[
\operatorname{PPL}=\exp\left(-\frac1T\sum_{t=1}^T\ln r_t\right)
=\left(\prod_{t=1}^T\frac1{r_t}\right)^{1/T}.
\]
La igualdad se obtiene aplicando propiedades del logaritmo. Si $r_t=0$
para algún evento observado, la pérdida y la perplejidad son infinitas.
Si todos los eventos reciben probabilidad $1/K$, la PPL vale $K$.
\subsubsection{Producto de probabilidades y probabilidad inicial}
Una matriz de transición por sí sola solo especifica las probabilidades
condicionales de pasos consecutivos. Para dar una ley sobre cadenas de
longitud $T$ se debe fijar $p(X_1)$ o condicionar explícitamente al símbolo
inicial. Si la prueba \texttt{abca} tiene transición total $1/35$ bajo
$\alpha=1$ y la inicial es uniforme, su probabilidad conjunta es $1/105$.
Este detalle afecta el número de factores $T$ y el denominador de la PPL.
\subsubsection{Bits por carácter y unidades de información}
Para el texto $z$ con $N$ caracteres bajo una convención documentada,
\[
\operatorname{BPC}(z;q)=\frac{-\log_2q(\tau(z))}{N}.
\]
Si $\tau$ es una tokenización reversible y la probabilidad cubre el texto
completo y todas las marcas que lo determinan, esta cantidad representa
una tasa ideal por carácter. De otro modo puede faltar información para
reconstruir el texto. El nombre BPC requiere especificar si se cuentan
puntos de código Unicode, bytes o grafemas y cómo se tratan saltos de línea.
\subsubsection{Relación entre PPL y BPC bajo una tokenización fijada}
Si la segmentación de $z$ contiene $T$ tokens y $\operatorname{PPL}_q$
se calculó sobre exactamente esos $T$ eventos, entonces
\[
\operatorname{BPC}=\frac{T}{N}\log_2\operatorname{PPL}_q.
\]
La fórmula es algebraica; no autoriza mezclar valores de PPL, $T$ y $N$
provenientes de convenciones o textos distintos.
\subsubsection{Byte Pair Encoding: reducción local}
Sea un corpus representado por secuencias finitas y un par $(u,v)$.
Cada reemplazo \emph{no solapado} de $u,v$ por $uv$ disminuye
el número de símbolos en una unidad. Si se ejecutan $r$ reemplazos,
la longitud total desciende exactamente $r$. Para el par $(a,a)$ en
$\texttt{aaaaa}$ se cuentan cuatro apariciones solapadas pero solo dos
reemplazos no solapados de izquierda a derecha: la disminución es dos.
Por ello, una implementación que cuenta pares brutos debe separar
``frecuencia'' de ``reemplazos efectivos''.
\subsubsection{Por qué BPE no prueba un óptimo global}
El algoritmo elige, en cada paso, el par de mayor frecuencia y toma una
decisión local. Aunque disminuye la longitud inmediatamente, cambiar las
fusiones tempranas cambia el conjunto de pares futuros; por ello la
propiedad local no implica optimalidad global con un presupuesto dado de
fusiones. También debe documentarse la política de desempate.
\subsubsection{Preguntas suplementarias de nivel matemático}
\begin{enumerate}
\item Pruebe la concavidad de $H(p)$ usando el Hessiano restringido al
interior del símplex.
\item Derive la regla de la cadena de KL a partir de las factorizaciones
$p(x_{1:T})$ y $q(x_{1:T})$.
\item Estudie el límite $\alpha\to\infty$ del estimador suavizado
con $K$ fijo, e interprete el resultado.
\item Dé un ejemplo donde dos modelos tengan igual pérdida en bits por
carácter y perplejidades por token distintas.
\item Explique por qué la validación aleatoria por caracteres es
incorrecta cuando el corpus contiene documentos muy similares.
\end{enumerate}
\clearpage

% ===== Fuente fuentes/propiedades_aplicadas.tex =====
\subsection{Dos propiedades que conectan información y modelado}


\subsubsection*{La información no aumenta al agrupar categorías}

Supongamos que un experimento observa tres categorías, pero un instrumento
solo informa si ocurrió $a$ o si ocurrió alguna de las otras dos.
La transformación $f: \{a,b,c\}\to\{a,\mathrm{otro}\}$ agrupa
los resultados $b$ y $c$. Para
$p=(1/2,1/4,1/4)$ y $q=(1/4,1/2,1/4)$, teníamos
$D_{\rm KL}(p\|q)=1/4$ bit. Después de la agregación,
\[
 p_f=(1/2,1/2),\qquad q_f=(1/4,3/4),
\]
y la divergencia entre las distribuciones observables es
\[
\begin{split}
D_{\rm KL}(p_f\|q_f)
&=\frac12\log_2\frac{1/2}{1/4}
+\frac12\log_2\frac{1/2}{3/4}\\
&=\frac12+\frac12\log_2\frac23
\approx0.207519\ \text{bits}.
\end{split}
\]
Ha disminuido respecto del $0.25$ original. Este ejemplo particular
ilustra una propiedad general de la KL denominada \emph{procesamiento de
datos}: al aplicar una transformación determinista común a los
resultados, no aumenta la divergencia.

\begin{figure}[H]\centering
\begin{tikzpicture}[>=Stealth,font=\small,
 node/.style={draw=azul,rounded corners,fill=azulclaro,minimum width=17mm,minimum height=8mm}]
\node[node](a)at(0,1){$a$};\node[node](b)at(0,0){$b$};\node[node](c)at(0,-1){$c$};
\node[node](yes)at(6,0.8){$a$};\node[node](oth)at(6,-0.8){otro};
\draw[->] (a)--(yes);\draw[->] (b)--(oth);\draw[->] (c)--(oth);
\node at (3,-1.8) {La observación final olvida la distinción entre $b$ y $c$.};
\end{tikzpicture}
\caption{La agregación pierde información sobre la diferencia entre dos leyes.}
\end{figure}

\textbf{Justificación.} Si $A_y=f^{-1}(y)$, la desigualdad de suma logarítmica
permite escribir
\[
\sum_{i\in A_y}p_i\ln\frac{p_i}{q_i}
\ge\left(\sum_{i\in A_y}p_i\right)
\ln\frac{\sum_{i\in A_y}p_i}{\sum_{i\in A_y}q_i}.
\]
Sumando en las categorías agregadas se obtiene
$D_{\rm KL}(p\|q)\ge D_{\rm KL}(p_f\|q_f)$.
La prueba exige las convenciones habituales de soporte, y no afirma
que ambas divergencias tengan que ser diferentes: podrían coincidir.
Este resultado es útil en investigación aplicada cuando un reporte
técnico se reduce a etiquetas de eventos; esa reducción puede eliminar
diferencias que existían a nivel textual.

\subsubsection*{Por qué ampliar el contexto puede reducir la incertidumbre óptima}

Sea $X$ el próximo símbolo, $Y$ el último símbolo observado y $Z$
una porción anterior del contexto. La teoría de la información
establece
\[
H(X\mid Y,Z)\le H(X\mid Y).
\]
No es una afirmación acerca de que \emph{cualquier modelo entrenado}
con más contexto rendirá mejor. Es una desigualdad sobre las verdaderas
entropías condicionales de una distribución fija.

Para justificarla, definamos la información mutua condicional:
\[
I(X;Z\mid Y)
=\mathbb E_{Y,Z}\Bigl[D_{\rm KL}\bigl(
 p(X\mid Y,Z)\,\|\,p(X\mid Y)
\bigr)\Bigr].
\]
Por Gibbs, el integrando es no negativo en cada contexto del soporte.
Al expandir los logaritmos también se obtiene
\[
I(X;Z\mid Y)=H(X\mid Y)-H(X\mid Y,Z).
\]
Combinando ambas igualdades deducimos la desigualdad.
La igualdad ocurre si $X$ y $Z$ son condicionalmente independientes dado $Y$.

\begin{ejemplo}[Una dependencia que el último símbolo no explica]
En la construcción $X_1$ y $X_2$ Bernoulli$(1/2)$ independientes y
$X_3=X_1$, se tiene
\[
H(X_3\mid X_2)=1\ \text{bit},
\qquad H(X_3\mid X_1,X_2)=0.
\]
La historia más amplia elimina toda incertidumbre \emph{verdadera}
sobre $X_3$, mientras que el bigrama conserva un bit de incertidumbre.
En muestras pequeñas, sin embargo, estimar más contextos puede aumentar
la varianza y empeorar la predicción en prueba. Allí reaparece el
compromiso entre capacidad del modelo, datos y suavizado.
\end{ejemplo}
La transición de este resultado a los Transformers es conceptual:
el acceso a contexto amplio \emph{permite} representar dependencias
que los bigramas descartan, pero aprenderlas de forma fiable exige
datos, parametrización, optimización y evaluación, asuntos de S02--S05.

% ===== Fin fuente =====


% ===== Fuente fuentes/puentes_aplicados.tex =====
\subsection{Profundización: cinco conexiones que preparan la investigación}


\subsubsection*{1. ¿Qué se puede concluir de una cadena de Markov?}

El modelo suavizado del ejemplo principal puede estudiarse como una cadena
homogénea de tres estados, con distribución inicial $\pi_0$ y matriz
\[
P=\begin{pmatrix}1/7&3/7&3/7\\1/2&1/3&1/6\\2/5&2/5&1/5\end{pmatrix}.
\]
Por construcción, $P_{ij}\ge0$ y $\sum_jP_{ij}=1$; cada fila define
una probabilidad. Como todas las entradas son estrictamente positivas,
la cadena es irreducible y aperiódica: desde cualquier estado se llega
a cualquier otro en un paso, y cada estado admite una transición a sí mismo.
Bajo estas condiciones existe una única distribución estacionaria $\pi$
con componentes estrictamente positivas y $\pi P=\pi$.

Para calcularla, escribimos $\pi=(u,v,w)$ y resolvemos
\[
\begin{cases}
 u=\tfrac17u+\tfrac12v+\tfrac25w,\\
 v=\tfrac37u+\tfrac13v+\tfrac25w,\\
 w=\tfrac37u+\tfrac16v+\tfrac15w,\\
 u+v+w=1.
\end{cases}
\]
Una de las primeras tres ecuaciones es redundante: las filas de $P$
suman uno. Tras sustituir y resolver el sistema lineal queda
\[
\boxed{\pi=\left(\frac{98}{281},\frac{108}{281},\frac{75}{281}\right).}
\]
Se verifica, por ejemplo, que
\[
\frac{98}{281}\frac17+\frac{108}{281}\frac12
+\frac{75}{281}\frac25
=\frac{14+54+30}{281}=\frac{98}{281}.
\]
El mismo procedimiento verifica las otras dos coordenadas.
Este cálculo revela la distinción entre \emph{probabilidad de transición}
y \emph{probabilidad de ocupación a largo plazo}: no son iguales por
definición. Una distribución inicial arbitraria evoluciona como
$\pi_t=\pi_0P^t$ y se aproxima a la estacionaria bajo las hipótesis
mencionadas. En lenguaje real, las probabilidades condicionadas suelen
cambiar con el contexto y no existe necesariamente una matriz homogénea
fija en ese sentido.

\begin{puente}
Este es el punto de contacto directo con el curso paralelo de Series
de Tiempo: cadenas de Markov, matrices estocásticas, distribución invariante
y pronóstico condicional. Más adelante podremos comparar esta representación
con modelos ocultos de Markov y modelos de estado, pero en S01 no es
necesario introducir sus algoritmos.
\end{puente}

\subsubsection*{2. Entropía condicional y tasa de información de una cadena}

La entropía conjunta de dos símbolos satisface
\[
H(X_1,X_2)=H(X_1)+H(X_2\mid X_1).
\]
Esto se demuestra al expandir
$-\sum_{i,j}p(i,j)\log_2 p(i,j)$ y usar
$p(i,j)=p(i)p(j\mid i)$:
\[
\begin{split}
H(X_1,X_2)&=-\sum_{i,j}p(i,j)\log_2p(i)\\
&\quad-\sum_{i,j}p(i,j)\log_2p(j\mid i)\\
&=H(X_1)+\sum_i p(i)H(X_2\mid X_1=i).
\end{split}
\]
El promedio $\sum_i p(i)H(X_2\mid X_1=i)$ es
la entropía condicional: mide cuánto falta por conocer sobre $X_2$
cuando se conoce el estado precedente.

Para la cadena estacionaria construida antes, su tasa de entropía es
\[
h(P,\pi)=-\sum_{i\in V}\pi_i\sum_{j\in V}P_{ij}\log_2P_{ij}.
\]
La fórmula surge porque el símbolo inicial contribuye un término finito
$H(X_1)$, mientras que cada transición, bajo estacionariedad, tiene
la misma entropía condicional media. En efecto,
\[
H(X_{1:T})=H(X_1)+\sum_{t=2}^T H(X_t\mid X_{t-1})
=H(\pi)+(T-1)h(P,\pi).
\]
Dividiendo por $T$ y dejando $T\to\infty$ se obtiene
$\lim_{T\to\infty}T^{-1}H(X_{1:T})=h(P,\pi)$.
Hay dos lecciones: el costo de representar el símbolo inicial
no desaparece en una cadena corta y la entropía marginal $H(\pi)$
no coincide necesariamente con la entropía condicional media.

\begin{ejemplo}[Un contraste con números sencillos]
Si las transiciones son independientes del estado anterior y cada fila
es $(1/2,1/2)$, entonces $H(X_t\mid X_{t-1})=1$ bit y la cadena
no aprovecha memoria. En cambio, para
$P=\left(\begin{smallmatrix}0.9&0.1\\0.1&0.9\end{smallmatrix}\right)$
la estacionaria es uniforme, $H(X_t)=1$ bit, pero
\[
H(X_t\mid X_{t-1})=-0.9\log_20.9-0.1\log_20.1\approx0.4690
\text{ bits}.
\]
El conocimiento del estado anterior reduce la incertidumbre de
la siguiente observación sin cambiar la distribución marginal.
\end{ejemplo}

\subsubsection*{3. Soporte, símbolos desconocidos y final de secuencia}

El suavizado aditivo distribuye masa positiva entre los $K$ símbolos
incluidos en el vocabulario del modelo. \emph{No asigna automáticamente}
probabilidad a cualquier símbolo que no pertenezca a ese vocabulario.
En un modelo de palabras entrenado con $V$, una palabra nueva necesita
una decisión de representación: sustituirse por un símbolo
$\langle\mathrm{UNK}\rangle$, segmentarse en unidades menores o
admitirse mediante algún mecanismo adicional.

\begin{figure}[H]\centering
\begin{tikzpicture}[>=Stealth,font=\small,
 box/.style={draw=azul,fill=azulclaro,rounded corners,align=center,text width=3.1cm,minimum height=9mm}]
\node[box](a)at(0,0){Palabra de prueba\\\texttt{revestidor}};
\node[box](b)at(4.1,1){Vocabulario de palabras\\\texttt{<UNK>}};
\node[box](c)at(4.1,-1){Vocabulario subléxico\\varias unidades};
\draw[->,thick] (a)--(b);\draw[->,thick] (a)--(c);
\end{tikzpicture}
\caption{Dos políticas distintas para un símbolo no visto; no son intercambiables.}
\end{figure}

Si además la longitud de la cadena es variable, necesitamos fijar
cómo termina la secuencia. Una construcción usual introduce un símbolo
de fin $\langle/\mathrm{s}\rangle$ y asigna probabilidades a las
transiciones hasta encontrarlo. De este modo, bajo hipótesis de
terminación adecuadas, el modelo determina una ley sobre cadenas
finitas de longitudes variables. Si se omite la probabilidad de parar,
la suma de probabilidades de cadenas completas puede no estar
bien definida.

\begin{ejemplo}[Parada con probabilidad geométrica]
Considérese un generador que emite un símbolo ordinario con probabilidad
$1-r$ y termina con probabilidad $r\in(0,1)$ en cada paso, con decisiones
independientes. Si $N$ es el número de símbolos ordinarios previos al fin,
\[
P(N=n)=(1-r)^nr,\quad n=0,1,2,\ldots,
\qquad \mathbb E[N]=\frac{1-r}{r}.
\]
Al limitar la generación a un máximo de $m$ símbolos mediante un arnés,
la probabilidad de truncamiento es $P(N\ge m)=(1-r)^m$.
Esta cuenta es una idealización explícita: en un LLM real la probabilidad
de fin depende de la historia y la política de decodificación.
\end{ejemplo}

\subsubsection*{4. Códigos ideales y límites de interpretación}

La cantidad $-\log_2q(x)$ se interpreta como longitud de código ideal
según $q$. Esto motiva tanto la entropía como la divergencia KL:
si los datos siguen $p$ pero se utiliza $q$,
\[
\mathbb E_p[-\log_2q(X)]-\mathbb E_p[-\log_2p(X)]
=D_{\rm KL}(p\|q).
\]
Es el exceso esperado de longitud ideal, no una promesa de que
cualquier compresor real tenga exactamente esa longitud en un archivo.
Los códigos binarios de longitud entera deben cumplir restricciones
como la desigualdad de Kraft y agregan sobrecostos de cabecera,
fronteras y representación.

\begin{ejemplo}[Cuando la entropía no es el costo de un símbolo concreto]
Para $p=(1/2,1/4,1/4)$, el símbolo de probabilidad $1/2$ contribuye
$1$ bit si ocurre, mientras que cada símbolo de probabilidad $1/4$
contribuye $2$ bits. La entropía $1.5$ bits es un \emph{promedio esperado},
no la longitud de cada realización. Con un código de prefijo
$\{0,10,11\}$ las longitudes individuales son exactamente $(1,2,2)$,
y el promedio coincide con 1.5 bits.
\end{ejemplo}

La pérdida logarítmica es una regla de puntuación propia: su riesgo
esperado se minimiza al informar la distribución verdadera.
La demostración ya está contenida en la identidad de entropía cruzada
más Gibbs. Por ello tiene sentido usarla para evaluar predicciones
probabilísticas, pero no sustituye métricas de exactitud, fidelidad o
utilidad operacional cuando esas son las preguntas científicas.

\subsubsection*{5. Incertidumbre experimental y unidad de remuestreo}

Supongamos que se quiere comparar dos procedimientos A y B en $m$
documentos reservados. Si $L_{A,j}$ y $L_{B,j}$ denotan las pérdidas
en el mismo documento $j$, definimos
\[
d_j=L_{A,j}-L_{B,j},\qquad \overline d=\frac1m\sum_{j=1}^m d_j.
\]
Un análisis pareado retiene la estructura de comparación: cada documento
aporta una diferencia. Si las unidades documentales son razonablemente
independientes, su error estándar clásico es $s_d/\sqrt m$, donde
\[
s_d^2=\frac1{m-1}\sum_{j=1}^m(d_j-\overline d)^2.
\]
Este error estándar y los intervalos que se construyen con él no
resultan válidos por el mero hecho de tener millones de tokens: los
símbolos dentro de un mismo documento comparten dependencias.

En investigación con textos técnicos, varios documentos del mismo pozo,
incidente o empresa también pueden estar correlacionados. En ese caso la
unidad de remuestreo puede necesitar elevarse hasta el pozo, incidente o
fuente, según la pregunta y el proceso de muestreo. Una comparación que
ignora esto puede aparentar una precisión mayor de la real.

\begin{ejemplo}[Diferencias pareadas construidas para ilustrar el método]
Considérense cuatro documentos con diferencias en BPC
$(-0.20,0.10,-0.10,0.00)$, calculadas como A menos B.
Su promedio es $-0.05$ BPC. Los cuadrados respecto de la media son
$0.0225,0.0225,0.0025,0.0025$; la suma es $0.05$.
La varianza muestral es $0.05/3\approx0.01667$ y el error estándar
$\sqrt{0.01667/4}\approx0.06455$. Con solo cuatro documentos,
el valor puntual no debe venderse como efecto establecido. El ejemplo es
\textbf{ILUSTRATIVO}, no procede de una prueba real.
\end{ejemplo}

\subsubsection*{6. Cuidado con una supuesta ventaja metodológica}

Un tokenizador puede producir menos unidades por palabra sin que
su modelo probabilístico correspondiente alcance menor pérdida de
predicción; la arquitectura, la cantidad de datos y el tratamiento
de caracteres pueden alterar esa relación. Por eso conviene separar
métricas \emph{descriptivas de segmentación} de métricas \emph{predictivas}.
Si comparamos dos modelos completos, podremos describir diferencias
de tasa; si queremos atribuir la diferencia a BPE, habrá que diseñar un
experimento en el que los otros factores se mantengan comparables.

Esta cautela también explica por qué la matriz IC10 es una \emph{semilla}
de investigación, no un resultado preestablecido. Un resultado nulo,
negativo o contrario a la hipótesis puede ser científicamente útil si
la medición fue defendible. El paso siguiente, en S02, será revisar
candidatos y disponibilidad de datos antes de elegir un problema.

% ===== Fin fuente =====


% ===== Fin fuente =====


% ===== Fuente modulos/70_ejercicios_contenido.tex =====
\phantomsection\label{s01:mod:ej}

% ===== Fuente fuentes/ejercicios_guia.tex =====
\subsection{Ejercicios matemáticos de la guía}

\noindent Estos ejercicios consolidan resultados de la clase y corresponden al
trabajo autónomo, no al tiempo presencial. Las soluciones justificadas están
separadas en \texttt{solej01.tex}.
\begin{enumerate}
\item \textbf{[Derivación]} Demuestre la regla de la cadena para $T=4$.
Para $X_1,X_2$ Bernoulli$(1/2)$ independientes y $X_3=X_1$,
contraste $p(1,0,1)$ con la aproximación de orden uno.
\emph{Autoverificación:} $1/4$ frente a $1/8$.
\item \textbf{[Cálculo]} Para $c=(6,3,1,0,0)$ y $K=5$,
calcule la predictiva usando dos \emph{priors} simétricas:
Dirichlet con parámetros unitarios y Dirichlet con parámetros $1/2$.
Compare la masa sobre categorías no vistas.
\emph{Autoverificación:} $1/15$ frente a $1/25$ en cada categoría nula.
\item \textbf{[Demostración]} Pruebe Gibbs desde $\ln t\le t-1$,
justifique cuándo se da la igualdad y deduzca
$H(p,q)=H(p)+D_{\rm KL}(p\|q)$.
\emph{Autoverificación:} KL es no negativa, pero generalmente asimétrica.
\item \textbf{[Cálculo]} Un predictor asigna a cuatro resultados observados
$(1/2,1/4,1/8,1/2)$. Calcule pérdida total, pérdida media y
PPL. Repita cuando la tercera probabilidad sea $1/16$.
\emph{Autoverificación:} PPL inicial $2^{7/4}$, final $4$.
\item \textbf{[Aplicación]} En un mismo texto de 100 caracteres, A
produce 100 tokens con PPL 3.2 y B 30 tokens con PPL 25.
Convierta a bits totales y BPC y enumere cuatro condiciones necesarias
para interpretar una diferencia.
\emph{Autoverificación:} BPC(A) $\approx1.6781$;
BPC(B) $\approx1.3932$.
\end{enumerate}
\subsubsection*{Problemas suplementarios}
\begin{enumerate}
\item[(S1.A)] Para una cadena de Markov homogénea, relacione su
entropía condicional $H(X_t\mid X_{t-1})$ con las tasas de
codificación cuando existe una distribución estacionaria.
\item[(S1.B)] Formalice por qué una frecuencia bruta de pares solapados
no es la reducción de longitud de una fusión BPE.
\item[(S1.C)] \emph{Semilla IC10:} formule una hipótesis falsable que
compare tasa de segmentación en corpus general y técnico, e indique
qué control haría falta para atribuir un efecto al tokenizador.
\end{enumerate}

% ===== Fin fuente =====


% ===== Fin fuente =====


% ===== Fuente fuentes/bibliografia_revisada.tex =====

\subsection{Lecturas y bibliografía de la sesión}
\phantomsection
\textbf{Entrada pedagógica:} el capítulo 6 de la primera edición de
\emph{Deep Learning with R} ofrece una transición visual entre texto,
tokens y representaciones; las ediciones segunda y tercera conectan esas
ideas con modelos neuronales posteriores.

\textbf{Fundamentos matemáticos:} entropía y codificación en Shannon;
propiedades de entropía, KL y modelado probabilístico en Cover y Thomas;
introducción a n-gramas y tokenización en Jurafsky y Martin.

\textbf{Algoritmos y ampliaciones:} BPE de subpalabras en Sennrich y
colaboradores; suavizado avanzado en Kneser y Ney. La implementación BPE
de S01 es una adaptación didactica y no reproduce un tokenizador industrial.


% ===== Fin fuente =====

